\documentclass[letterpaper,11pt,onesided]{book}%

\newif\ifinstructor
%% uncomment one of the following, depending if you want to include the
%% instructor notes or not
\instructortrue
%\instructorfalse


\newif\ifnotes
%% uncomment one of the following, depending if you want to include the
%% instructor notes or not
\notestrue
%\notesfalse

\usepackage{tabls}
\usepackage{booktabs}
\usepackage{amsmath}
\usepackage{amssymb}
\usepackage{amsthm}
\usepackage{amsfonts}
\usepackage{multicol}
\usepackage{enumitem}
\usepackage{stmaryrd}
\usepackage{wasysym}
\usepackage{microtype}
\usepackage{tikz}
\usetikzlibrary{positioning}
\usepackage[hang,flushmargin,symbol*]{footmisc}
\usepackage{comment}
\usepackage{graphicx}
\usepackage{wrapfig}
\usepackage{color}
\definecolor{darkblue}{rgb}{0, 0, .6}
\definecolor{grey}{rgb}{.7, .7, .7}
\usepackage[breaklinks]{hyperref}
\hypersetup{
	colorlinks=true,
	linkcolor=darkblue,
	anchorcolor=darkblue,
	citecolor=darkblue,
	pagecolor=darkblue,
	urlcolor=darkblue,
	pdftitle={},
	pdfauthor={}
}

\let\oldsection\section
\renewcommand\section{\clearpage\oldsection}


\usepackage[margin=1in]{geometry}

%%%%%%%%%%%%%%%%%%%%%%%%%%%%%%%%%%%%%%%%%%%%%%%%%%%%%%%%%%%%%%%%%%%%%%%%%%%%%%%%%%%%%%%%%%%%%%%%%%%%%%
% Margin notes
%%%%%%%%%%%%%%%%%%%%%%%%%%%%%%%%%%%%%%%%%%%%%%%%%%%%%%%%%%%%%%%%%%%%%%%%%%%%%%%%%%%%%%%%%%%%%%%%%%%%%%
%\usepackage[left=1in,right=2.75in,top=1in,bottom=1in]{geometry}
%\marginparwidth 1.75in
%\let\oldmarginpar\marginpar
%\renewcommand\marginpar[1]{\-\oldmarginpar{\raggedright\footnotesize #1}}
%\renewcommand\marginpar[1]{\-\oldmarginpar[\raggedleft\footnotesize #1]{\raggedright\footnotesize #1}}




%%%%%%%%%%%%%%%%%%%%%%%%%%%%%%%%%%%%%%%%%%%%%%%%%%%%%%%%%%%%%%%%%%%%%%%%%%%%%%%%%%%%%%%%%%%%%%%%%%%%%%
% Embed the date and time into every page
%%%%%%%%%%%%%%%%%%%%%%%%%%%%%%%%%%%%%%%%%%%%%%%%%%%%%%%%%%%%%%%%%%%%%%%%%%%%%%%%%%%%%%%%%%%%%%%%%%%%%%
%\usepackage[12hr]{datetime}
%\newdateformat{draftdate}{%
%\shortdayofweekname{\THEDAY}{\THEMONTH}{\THEYEAR}, %
%\THEDAY\ \shortmonthname[\THEMONTH] \THEYEAR}
%\draftdate
%\usepackage{eso-pic}
%\AddToShipoutPicture{\put(10,10){\small Draft: \today\ at \currenttime }}%--- version: \MakeUppercase{\svnInfoRevision}}}

%%%%%%%%%%%%%%%%%%%%%%%%%%%%%%%%%%%%%%%%%%%%%%%%%%%%%%%%%%%%%%%%%%%%%%%%%%%%%%%%%%%%%%%%%%%%%%%%%%%%%%
% Instructor-specific material (answers, helps, etc.)
%%%%%%%%%%%%%%%%%%%%%%%%%%%%%%%%%%%%%%%%%%%%%%%%%%%%%%%%%%%%%%%%%%%%%%%%%%%%%%%%%%%%%%%%%%%%%%%%%%%%%%
\ifinstructor
  \newcommand{\instructor}[1]{\marginpar{\textbf{Instructor: }#1}}
\else
  \newcommand{\instructor}[1]{}
\fi

%%%%%%%%%%%%%%%%%%%%%%%%%%%%%%%%%%%%%%%%%%%%%%%%%%%%%%%%%%%%%%%%%%%%%%%%%%%%%%%%%%%%%%%%%%%%%%%%%%%%%%
% Personal notes about improving the text, etc.
%%%%%%%%%%%%%%%%%%%%%%%%%%%%%%%%%%%%%%%%%%%%%%%%%%%%%%%%%%%%%%%%%%%%%%%%%%%%%%%%%%%%%%%%%%%%%%%%%%%%%%
\ifnotes
\renewcommand{\thefootnote}{\roman{footnote}}
\newcommand{\note}[1]{\footnote{#1}\marginpar{\fbox{\textbf{\thefootnote}}}}
\else
\newcommand{\note}[1]{}
\fi



%%%%%%%%%%%%%%%%%%%%%%%%%%%%%%%%%%%%%%%%%%%%%%%%%%%%%%%%%%%%%%%%%%%%%%%%%%%%%%%%%%%%%%%%%%%%%%%%%%%%%%
% Commands to define a term (so that we get nice indexing
%%%%%%%%%%%%%%%%%%%%%%%%%%%%%%%%%%%%%%%%%%%%%%%%%%%%%%%%%%%%%%%%%%%%%%%%%%%%%%%%%%%%%%%%%%%%%%%%%%%%%%
\newcommand{\define}[2][]{\textbf{#2}\index{#1#2}}
\newcommand{\definemargin}[2][]{\marginpar{#2}\define[#1]{#2}}

%%%%%%%%%%%%%%%%%%%%%%%%%%%%%%%%%%%%%%%%%%%%%%%%%%%%%%%%%%%%%%%%%%%%%%%%%%%%%%%%%%%%%%%%%%%%%%%%%%%%%%
% Theorems
%%%%%%%%%%%%%%%%%%%%%%%%%%%%%%%%%%%%%%%%%%%%%%%%%%%%%%%%%%%%%%%%%%%%%%%%%%%%%%%%%%%%%%%%%%%%%%%%%%%%%%

\newtheoremstyle{box}%
{}{}% standard spacing before and after
{}% Body style
{}{\bfseries}{.}% Heading indent, font, and punctuation
{ }% space after heading
{\thmname{#1}\thmnumber{ #2}\thmnote{: #3}}% head spec

\newtheoremstyle{problem}%
{}{}% standard spacing before and after
{}% Body style
{}{\bfseries}{}% Heading indent, font, and punctuation
{1em}% space after heading
{\fbox{\thmname{#1}\thmnumber{ #2}}\thmnote{ (#3)}}% head spec


\theoremstyle{problem}
\newtheorem{theorem}{Theorem}[chapter]
\newtheorem{question}[theorem]{Question}
\newtheorem{problem}[theorem]{Problem}
\newtheorem{exercise}[theorem]{Exercise}
\newtheorem{corollary}[theorem]{Corollary}

\theoremstyle{definition}
\newtheorem{acknowledgement}[theorem]{Acknowledgement}
\newtheorem{algorithm}[theorem]{Algorithm}
\newtheorem{axiom}[theorem]{Axiom}
\newtheorem{case}[theorem]{Case}
\newtheorem{claim}[theorem]{Claim}
\newtheorem{conclusion}[theorem]{Conclusion}
\newtheorem{condition}[theorem]{Condition}
\newtheorem{conjecture}[theorem]{Conjecture}
\newtheorem{criterion}[theorem]{Criterion}
\newtheorem{definition}[theorem]{Definition}
\newtheorem{example}[theorem]{Example}
\newtheorem{journal}[theorem]{Journal}
\newtheorem{lemma}[theorem]{Lemma}
\newtheorem{notation}[theorem]{Notation}
\newtheorem{proposition}[theorem]{Proposition}
\newtheorem{remark}[theorem]{Remark}
\newtheorem{solution}[theorem]{Solution}
\newtheorem{summary}[theorem]{Summary}
\newtheorem{skeleton}[theorem]{Skeleton Proof}


\newsavebox{\savepar}
\newenvironment{textbox}{\noindent\begin{lrbox}{\savepar}\begin{minipage}[c]{.98\textwidth}}{\end{minipage}\end{lrbox}\fcolorbox{black}{white}{\usebox{\savepar}}}


%%%%%%%%%%%%%%%%%%%%%%%%%%%%%%%%%%%%%%%%%%%%%%%%%%%%%%%%%%%%%%%%%%%%%%%%%%%%%%%%%%%%%%%%%%%%%%%%%%%%%%
% Notation shortcuts
%%%%%%%%%%%%%%%%%%%%%%%%%%%%%%%%%%%%%%%%%%%%%%%%%%%%%%%%%%%%%%%%%%%%%%%%%%%%%%%%%%%%%%%%%%%%%%%%%%%%%%

\DeclareMathOperator{\dom}{Dom}
\DeclareMathOperator{\codom}{Codom}
\DeclareMathOperator{\range}{Rng}
\newcommand{\blank}[1]{\raisebox{0pt}[14pt]{\rule{#1}{0.7pt}}}

\let\origdoublepage\cleardoublepage

\newcommand{\clearemptydoublepage}{%
  \clearpage
  {\pagestyle{empty}\origdoublepage}%
}



\begin{document}
\frontmatter
\title{Course Notes for Introduction to Proof}
\author{D.C.~Ernst\thanks{Northern Arizona University, \url{Dana.Ernst@nau.edu}}}
\date{Typeset on \today\\
\vfill
  Part of this work is an adaptation of notes written by Stan Yoshinobu of Cal Poly and Matthew Jones of California State University, Dominguez Hills.
\vfill
\includegraphics[height=1.3cm]{by-sa}
\vfill}
\maketitle
 \thispagestyle{empty}

% \noindent\copyright{ year author.  Some Rights Reserved.\\

% \bigskip

% \noindent Except where otherwise noted, this work is licensed under the Creative Commons
% Attribution-ShareAlike 3.0 United States License. To view a copy of
% this license, visit 
% \begin{center}
%   \url{http://creativecommons.org/licenses/by-sa/3.0/us/}
% \end{center}
% or send a letter to Creative Commons, 171 Second Street, Suite 300,
% San Francisco, California, 94105, USA.

% \bigskip

% \noindent Please attribute this work to:
% \smallskip

%  Dana Ernst (Northern Arizona University),  Stan Yoshinobu (Cal Poly), and Matthew Jones (California State University, Dominguez Hills)

% \vfill 

% The source for this work is available at \url{https://github.com/jasongrout/IBL-IntroToProof}.
% \vfill
% }
\tableofcontents
\clearemptydoublepage

\begin{chapter}{Preface}
This is preface text.
\end{chapter}

\clearemptydoublepage

\mainmatter

\begin{chapter}{Introduction to Mathematics}
Before you get started, make sure you've read Chapter 0, which sets the tone for the work we will begin doing here.

\begin{section}{A Taste of Number Theory}

%This section (or possibly a new section before this one) really needs more of a lead in.  For example, I should add:
%What is a theorem, proof, axiom, definition?
%What are the ground rules for doing presentations?
%Description of a direct proof.
%Discuss of implicit conditionals.

In this section, we will work with the set of integers, $\mathbb{Z} = \{\ldots, -3, -2, -1, 0, 1, 2, 3, \ldots\}$.  The purpose of this section is to get started with proving some theorems about numbers and study the properties of $\mathbb{Z}$.

It is important to note that we are diving in head first here.  There are going to be some subtle issues that you will bump into and our goal will be to see what those issues are, and then we will take a step back and start again.  See what you can do!

\begin{definition}
An integer $n$ is \textbf{even} if $n=2k$ for some integer $k$. \end{definition}

\begin{definition}
An integer $n$ is \textbf{odd} if $n=2k+1$ for some integer $k$. \end{definition}

\begin{theorem}\label{two consecutive ints}
The sum of two consecutive integers is odd.
\end{theorem}

\begin{theorem}
If $n$ is even, then $n^2$ is even.
\end{theorem}

\begin{problem}[*]\label{sum of even and odd}
Prove or provide a counterexample:  The sum of an even number and an odd number is odd.
\end{problem}

\begin{question}
Did Theorem~\ref{two consecutive ints} need to come before Problem~\ref{sum of even and odd}?  Could we have used Problem~\ref{sum of even and odd} to prove Theorem~\ref{two consecutive ints}?  If so, outline how this alternate proof would go.  Perhaps your original proof utilized the approach I'm hinting at.  If this is true, can you think of a proof that does not rely directly on Problem~\ref{sum of even and odd}?  Is one approach better than the other?
\end{question}

\begin{problem}
Prove or provide a counterexample: The product of an odd number and an even number is odd.
\end{problem}

\begin{problem}[*]
Prove or provide a counterexample: The product of an odd number and an odd number is odd.
\end{problem}

\begin{problem}[*]
Prove or provide a counterexample: The product of two even numbers is even.
\end{problem}

\begin{definition}
An integer $n$ \textbf{divides} the integer $m$, written $n|m$, if and only if there exists an integer $k$ such that $m=nk$. In the same context, we may also write that $m$ \textbf{is divisible by} $n$.
\end{definition}

In this section on number theory, we allow addition, subtraction, and multiplication.  Division is not allowed since an integer divided by an integer may result in a number that is not an integer. The upshot: don't write $\frac{m}{n}$.  When you feel the urge to divide, switch to an equivalent formulation using multiplication.

\begin{problem}
Let $n$ be an integer.  Prove or provide a counterexample: If 6 divides $n$, then 3 divides $n$.\end{problem}

\begin{problem}
Let $n$ be an integer.  Prove or provide a counterexample: If 6 divides $n$, then 4 divides $n$.
\end{problem}

\begin{theorem}[*]
Assume $n$, $m$, and $a$ are integers.  If $a|n$, then $a|mn$.
\end{theorem}

A theorem that follows almost immediately from another theorem is called a \textbf{corollary}.  See if you can prove the next result quickly using the previous theorem.  Be sure to cite the theorem in your proof.

\begin{corollary}
Assume $n$ and $a$ are integers.  If $a$ divides $n$, then $a$ divides $n^2$.
\end{corollary}

\begin{problem}
Assume $n$ and $a$ are integers.  Prove or provide a counterexample:  If $a$ divides $n^2$, then $a$ divides $n$.
\end{problem}

\begin{theorem}
Assume $a$ and $n$ are integers. If $a$ divides $n$, then $a$ divides $-n$. 
\end{theorem}

\begin{theorem}[*]
Assume $a$, $m$, and $n$ are integers. If $a$ divides $m$ and $a$ divides $n$, then $a$ divides $m+n$. 
\end{theorem}

Once we've proved a few theorems, we should be on the look out to see if we can utilize any of our current results to prove new results.  There's no point in reinventing the wheel if we don't have to.  Try to use a couple of our previous results to prove the next theorem.

\begin{theorem}
Assume $a$, $m$, and $n$ are integers. If $a$ divides $m$ and $a$ divides $n$, then $a$ divides $m-n$.
\end{theorem}

\begin{problem}[*]
Assume $a$, $b$, and $m$ are integers. Determine whether the following statement holds sometimes, always, or never.  If $ab$ divides $m$, then $a$ divides $m$ and $b$ divides $m$.  Justify with a proof or counterexample.
\end{problem}

\begin{theorem}[*]
If $a, b$, and $c$ are integers where $a$ divides $b$ and $b$ divides $c$, then $a$ divides $c$.
\end{theorem}

The previous theorem is referred to as \textbf{transitivity of division of integers}.

\begin{theorem}[*]
The sum of any three consecutive integers is always divisible by three.\end{theorem}

\end{section}

% unfortunately, this semester, there was an inconsistency in how
% theorems were numbered at the start of section 1.2.  Since we
% already passed out sheets for these, I am resetting the numbers
% here.
\addtocounter{theorem}{-1}
\begin{section}{Logic, Negation, Contrapositive}

After diving in head first in the last section, we'll take a step back and do a more careful examination of what it is we are actually doing.

\begin{definition}
A \textbf{proposition} (or \textbf{statement}) is a sentence that is either true or false.
\end{definition}

For example, the sentence ``All liberals are hippies'' is a false proposition.  However, the perfectly good sentence ``$x=1$'' is \emph{not} a proposition all by itself since we don't actually know what $x$ is.

\begin{exercise} Determine whether the following are propositions or not. Explain.
\begin{enumerate}
\item All cars are red.
\item Van Gogh was the best artist ever. 
\item If my name is Joe, then my name starts with the letter J.
\item If my name starts with the letter J, then my name is Joe.
\item $f$ is continuous.
\item All functions are continuous.
\item If $f$ is a differentiable function, then $f$ is continuous function.
\item The president had eggs for breakfast the morning of his tenth birthday.
\item What time is it? 
\item There exists an $x$ such that $x^2=4$.
\item $x^2=4.$
\item $\sqrt{2}$ is an irrational number.
\item For all real numbers $x$, $x^3=x$.
\item There exists a real number $x$ such that $x^3=x$.
\item $p$ is prime.
\end{enumerate}
\end{exercise}

Given two propositions, we can form more complicated propositions using the logical connectives ``and", ``or", and ``If\ldots, then\ldots".

\begin{definition}
Let $A$ and $B$ be propositions.  The proposition ``\textbf{$A$ and $B$}'' is true if and only if both $A$ and $B$ are true.  The statement ``$A$ and $B$'' is expressed symbolically as 
\[
A \wedge B
\]
and is known as the \textbf{conjunction} of $A$ and $B$.
\end{definition}

\begin{definition}
Let $A$ and $B$ be propositions.  The proposition ``\textbf{$A$ or $B$}'' is true if and only if at least one of $A$ or $B$ is true.  The statement ``$A$ or $B$'' is symbolically represented as
\[
A \vee B
\]
and is known as the \textbf{disjunction} of $A$ and $B$.
\end{definition}

\begin{definition}
Let $A$ be a proposition.  The \textbf{negation} of $A$, denoted $\neg A$, is true if and only if $A$ is false.
\end{definition}

\begin{exercise}Describe the meaning of $\neg (A \wedge B)$ and $\neg (A \vee B)$.
\end{exercise}

\begin{definition}
A \textbf{truth table} is a table that illustrates all possible truth values for a proposition.  
\end{definition}

\begin{example}
Let $A$ and $B$ be propositions.  Then the truth table for the conjunction $A\wedge B$ is given by the following.
\[
\begin{tabular}{| c | c  | c |} \hline
$A$  &  $B$ & $A \wedge B$  \\ \hline\hline
T & T & T  \\ 
T & F & F  \\ 
F & T & F  \\ 
F & F & F  \\ \hline
\end{tabular}
\]
Notice that we have columns for each of $A$ and $B$.  The rows for these two columns correspond to all possible combinations for $A$ and $B$.  The third column gives us the truth value of $A\wedge B$ given the possible truth values for $A$ and $B$.
\end{example}

\begin{exercise}
Create a truth table for each of $A \vee B$, $\neg A$, $\neg (A \wedge B)$, and $\neg A \wedge \neg B$.  Feel free to add additional columns to your tables to assist you with intermediate steps.
\end{exercise}

\begin{exercise}
Suppose $P$ is a complex proposition built out of the propositions $A$, $B$, and $C$.  How many rows would the truth table for $P$ require?
\end{exercise}

\begin{definition}\label{def:conditional}
Let $A$ and $B$ represent propositions.  The conditional proposition ``\textbf{If $A$, then $B$}'' is expressed symbolically as 
\[
A \implies B
\]
and has the following truth table.
\[
\begin{tabular}{| c | c  | c |} \hline
$A$  &  $B$ & $A \implies B$  \\ \hline\hline
T & T & T  \\ 
T & F & F  \\ 
F & T & T  \\ 
F & F & T  \\ \hline
\end{tabular}
\]
\end{definition}

\begin{exercise}\label{exer:translations}
Let $A$ represent ``6 is an even number'' and $B$ represent ``6 is a multiple of 4.''  Express each of the following in ordinary English sentences and state whether the statement is true or false.
\begin{enumerate}
  \item $A \wedge B$
  \item $A \vee B$
  \item $\neg A$
  \item $\neg B$
  \item $\neg (A \wedge B)$
  \item $\neg (A \vee B)$
  \item $A \implies B$
\end{enumerate}
\end{exercise}

\begin{problem}
Suppose I am the coach of our co-ed dodgeball team and you all are the players.  I tell you ``If we win tonight, then I will buy you pizza tomorrow.''  After reviewing the definition of conditional proposition, determine the case(s) in which you can rightly claim to have been lied to.
\end{problem}

\begin{definition}
Two statements are \textbf{logically equivalent} (or \textbf{equivalent} if the context is clear) if and only if they have the same truth table.  That is, proposition $P$ is true exactly when proposition $Q$ is true, and $P$ is false exactly when $Q$ is false.  When $P$ and $Q$ are logically equivalent we denote this symbolically as 
\[
P \iff Q,
\]
which we read ``$P$ if and only if $Q$".  It is common to abbreviate ``if and only if'' as ``iff''.
\end{definition}

Each of the next three facts can be justified using truth tables.

\begin{theorem}
If $A$ is a proposition, then $\neg(\neg A)$ is equivalent to $A$.
\end{theorem}

\begin{theorem}\label{thm:demorgan}
If $A$ and $B$ are propositions, then $\neg(A \wedge B) \iff \neg A \vee \neg B$.  (\emph{Note}: This theorem is referred to as DeMorgan's Law.)
\end{theorem}

\begin{problem}[*]
Let $A$ and $B$ be propositions.  Conjecture a statement similar to Theorem \ref{thm:demorgan} for the proposition $\neg(A\vee B)$ and then prove it.
\end{problem}

\begin{definition}
The \textbf{converse} of $A \implies B$ is $B \implies A$.
\end{definition}

\begin{definition}
The \textbf{contrapositive} of $A \implies B$ is $\neg B \implies \neg A$.
\end{definition}

\begin{exercise}
Let $A$ and $B$ represent the statements from Exercise \ref{exer:translations}.  Express the following in ordinary English sentences.
\begin{enumerate}
\item The converse of $A \implies B$
\item The contrapositive of $A \implies B$
\end{enumerate}
\end{exercise}

\begin{exercise} Find the contrapositive of the following statements: 
\begin{enumerate}
\item If $n$ is an even natural number, then $n+1$ is an odd natural number.\footnote{Despite the fact that each of ``$n$ is an even natural number" and ``$n+1$ is an odd natural number" are not propositions (since we cannot determine their truth values without knowing what $n$ is), the implication is a proposition.  We discuss this further when we introduce predicates.}
\item If it rains today, then I will bring my umbrella.
\item If it does not rain today, then I will not bring my umbrella.
\end{enumerate} \end{exercise}

\begin{exercise}
Provide an example of a true conditional proposition whose converse is false.
\end{exercise}

\begin{theorem}\label{thm:contrapos}
Assume $A$ and B are statements.  Then ${A\implies B}$ is equivalent to its contrapositive.
\end{theorem}

The upshot of Theorem \ref{thm:contrapos} is that if you want to prove a conditional proposition, you can prove its contrapositive instead.  Prove each of the next two propositions using the contrapositive of the given statement.

\begin{theorem}[*]
Assume $x$ and $y$ are integers.  If $xy$ is odd, then both $x$ and $y$ are odd. (Prove using contrapositive.)
\end{theorem}

\begin{theorem}[*]
Assume $x$ and $y$ are integers.  If $xy$ is even, then $x$ or $y$ is even.  (Prove using contrapositive.)
\end{theorem}

\end{section}

\begin{section}{Negating Implications and Proof by Contradiction}

So far we have discussed how to negate propositions of the form $A$, $A\wedge B$, and $A\vee B$ for propositions $A$ and $B$.  However, we have yet to discuss how to negate propositions of the form $A\implies B$.

\begin{problem}\label{prob:implication as disjunction}
Let $A$ and $B$ be propositions.  Conjecture an equivalent way of expressing the conditional proposition $A\implies B$ as a proposition involving the disjunction symbol $\vee$ and possibly the negation symbol $\neg$, but not the implication symbol $\implies$.  Prove your conjecture using a truth table.
\end{problem}

\begin{exercise}\label{exer:Darth Vader}
Let $A$ and $B$ be the propositions ``Darth Vader is a hippie" and ``Sarah Palin is a liberal", respectively.  Using Problem \ref{prob:implication as disjunction}, express $A\implies B$ as an English sentence involving the disjunction ``or."
\end{exercise}

\begin{problem}[*]\label{prob:negation of implication}
Let $A$ and $B$ be two propositions.  Conjecture an equivalent way of expressing the proposition $\neg(A\implies B)$ as a proposition involving the conjunction symbol $\wedge$ and possibly the negation symbol $\neg$, but not the implication symbol $\implies$.  Prove your conjecture using previous results.
\end{problem}

\begin{exercise}
Let $A$ and $B$ be the propositions in Exercise \ref{exer:Darth Vader}.  Using Problem \ref{prob:negation of implication}, express $\neg(A\implies B)$ as an English sentence involving the conjunction ``and."
\end{exercise}

\begin{exercise}
The following proposition is \emph{false}.  Negate this proposition to obtain a true statement.  Write your statement as a conjunction.
\begin{quote}
If $.\overline{99}=\frac{9}{10}+\frac{9}{100}+\frac{9}{1000}+\cdots$, then $.\overline{99}\neq 1$.
\end{quote}
You do \emph{not} need to prove your new statement.
\end{exercise}

Recall that a proposition is exclusively either true or false.  That is, a proposition can never be both true and false.  This idea leads us to the next definition.

\begin{definition}
A compound proposition that is always false is called a \textbf{contradiction}.  A compound statement that is always true is called a \textbf{tautology}.
\end{definition}

\begin{theorem}
Let $A$ be a proposition. Then $\neg A\wedge A$ is a contradiction.
\end{theorem}

\begin{exercise}
Provide an example of a tautology using arbitrary positions and any of the logical connectives $\neg$, $\wedge$, and $\vee$.  Then prove that your example is in fact a tautology.
\end{exercise}

Suppose that we want to prove some proposition $P$ (which might be something like $A\implies B$ or possibly more complicated).  One approach, called \textbf{proof by contradiction}, involves assuming $\neg P$ and then logically deducing a contradiction of the form $Q\wedge \neg Q$, where $Q$ is some proposition (possibly equal to $P$).  Since this is absurd, it cannot be the case that $\neg P$ is true, which implies that $P$ is true.  The tricky part about a proof by contradiction is that it is not usually obvious what the statement $Q$ is going to be.  Here is what the general structure for a proof by contradiction looks like.

\bigskip

\begin{skeleton}[Proof of $P$ by contradiction]
Here is what the general structure for a proof by contradiction looks like if we are trying to prove the proposition $P$.

\bigskip

\begin{textbox}
\begin{proof}
For sake of a contradiction, assume $\neg P$.
\begin{center}
$\vdots$\\
(Use definitions and previous theorems to derive some $Q$ and its negation $\neg Q$.)\\
$\vdots$
\end{center}
\noindent This is a contradiction.  Therefore, $P$.
\end{proof}
\end{textbox}

\end{skeleton}

Among other situations, proof by contradiction can be useful for proving statements of the form $A\implies B$, where $B$ is worded negatively or $\neg B$ is easier to ``get your hands on."  

\begin{skeleton}[Proof of $A\implies B$ by contradiction]\label{pf by contradiction for implication}
If you want to prove the proposition $A\implies B$ via a proof by contradiction, then the structure of the proof is as follows.

\bigskip

\begin{textbox}
\begin{proof}
For sake of a contradiction, assume $A$ and $\neg B$.
\begin{center}
$\vdots$\\
(Use definitions and previous theorems to derive some $Q$ and its negation $\neg Q$.)\\
$\vdots$
\end{center}
\noindent This is a contradiction.  Therefore, if $A$, then $B$.
\end{proof}
\end{textbox}
\end{skeleton}

\begin{question}
In Skeleton Proof \ref{pf by contradiction for implication}, why did we start by assuming $A$ and $\neg B$?
\end{question}

Prove the following theorem in two ways: (i) prove the contrapositive, and (ii) prove using a proof by contradiction.

\begin{theorem}[*]
Assume that $x\in\mathbb{Z}$.  If $x$ is odd, then 2 does not divide $x$. (Prove in two different ways.)
\end{theorem}

Prove the following theorem by contradiction.

\begin{theorem}[*]
Assume that $x,y\in\mathbb{N}$\footnote{Note that $\mathbb{N}=\{1,2,3,\ldots\}$ is the set of \textbf{natural numbers}. The given statement is not true if we replace $\mathbb{N}$ with $\mathbb{Z}$.  Do you see why?}.  If $x$ divides $y$, then $x\leq y$. (Prove using a proof by contradiction.)
\end{theorem}

\begin{question}
What obstacles (if any) are there to proving the previous theorem directly without using proof by contradiction?
\end{question}

\end{section}

% here, we skip 1 to get back on track (see above)...
\addtocounter{theorem}{1}
\begin{section}{Introduction to Quantification}

Recall that sentences of the form ``$x>0$" are not propositions (unless the context of $x$ is perfectly clear).  In this case, we call $x$ a \textbf{free variable}.  In order to turn a sentence with free variables into a proposition, for each free variable, we need to either substitute in a value (not necessarily a number) for the free variable or we must ``quantify" the free variable.  

\begin{definition}
A sentence with a free variable is called a \textbf{predicate}.
\end{definition}

\begin{exercise}
Give 3 examples of mathematical predicates involving 1, 2, and 3 free variables, respectively.
\end{exercise}

It is convenient to borrow function notation to represent predicates.  For example, each of the following represents a predicate with the indicated free variables.
\begin{itemize}
\item $S(x):=``x^2-4=0"$
\item $L(a,b):=``a<b"$
\item $F(x,y):=``x \mbox{ is friends with } y"$
\end{itemize}

The notation ``$:=$'' is often used in mathematics to show that we are defining something (as opposed to stating something is mathematically equal).  Also, note that the use of the quotation marks above removed some ambiguity.  If we wrote $S(x)=x^2-4=0$, it would look like $S(x)$ equals 0, but actually we want $S(x)$ to represent the whole sentence ``$x^2-4=0$''.

One way we can make propositions out of predicates is by assigning specific values to the free variables.  That is, if $P(x)$ is a predicate and $x_0$ is specific value, then $P(x_0)$ is now a proposition (and may be true or false).

\begin{exercise}
Consider $S(x)$ and $L(a,b)$ from the discussion above.  Determine the truth values of $S(2)$, $S(0)$, $S(-2)$, $L(1,2)$, $L(2,1)$, and $L(-3,-2)$.
\end{exercise}

\begin{exercise}
Again, consider $L(a,b)$ from above.  Is $L(2,b)$ a proposition or a predicate?  Explain your answer.
\end{exercise}

Besides substituting specific values in for free variables in a predicate, there are (at least) two other ways in which we can use predicates to build propositions.  We do this by making a claim about which values of the free variables apply to the predicate.

\begin{example}\label{ex:quantified predicates}
The following sentences are both propositions.
\begin{enumerate}
\item For all $x\in\mathbb{R}$, $x^2-4=0$.
\item There exists $x\in\mathbb{R}$ such that $x^2-4=0$.
\end{enumerate}
\end{example}

\begin{exercise}
Determine the truth value of the two propositions from Example \ref{ex:quantified predicates}.  What would it take to prove your answers?
\end{exercise}

\begin{definition}
``For all" is called the \textbf{universal quantifier} and ``there exists\ldots such that" is called the \textbf{existential quantifier}.
\end{definition}

The variables in propositions associated with quantifiers are called \textbf{bound variables}.  In order for a sentence containing variables to be a proposition, \emph{all} variables must be bound.  In other words, to be a proposition, all variables need to be quantified so that there is no ambiguity.

It is important to note that the existential quantifier is making a claim about ``at least one" \emph{not} ``exactly one."    Also, we can replace ``there exists\ldots such that" with phrases (possibly with some other tweaking to the sentence) like ``for some".  It is also worth noting that ``for all", ``for any", ``for every" are used interchangeably in mathematics (even though they might convey slightly different meanings in colloquial language).  

We must also take care that the universe of acceptable values for bound variables is clear.  Consider for a moment the proposition ``For all $x$, $x>0$".  Is this proposition true or false?  The answer depends on what $x$'s we are taking \emph{all} of.  For example, if the universe of discourse is the set of integers, then the proposition is false.  However, if we take the universe of acceptable values to be the natural numbers, then the proposition is true.  We must be careful to avoid such ambiguities.  Often, the context can resolve such ambiguities, but otherwise, we need to write things like: ``For all $x\in\mathbb{Z}$, $x>0$" or ``For all $x\in\mathbb{N}$, $x>0$".

\begin{exercise}
Suppose our universe of acceptable values is the set of integers.
\begin{enumerate}
\item Provide an example of a predicate $P(x)$ such that ``For all $x$, $P(x)$" is true.
\item Provide an example of a predicate $Q(x)$ such that ``For all $x$, $Q(x)$" is false, but ``There exists $x$ such that $Q(x)$" is true.
\end{enumerate}
\end{exercise}

If a predicate has more than one free variable, then we can build propositions by quantifying each variable.  However, the order of the quantifiers is extremely important!!!

\begin{exercise}\label{exer:ways to quantify}
Let $P(x,y)$ be a predicate with the free variables $x$ and $y$ (and let's assume the universe of discourse is clear).  Write down all possible ways (where order matters) that the variables could be quantified.  To get you started, here's one: For all $x$, there exists $y$ such that $P(x,y)$.  Find the rest.
\end{exercise}

\begin{problem}
Are there any propositions on your list from Exercise~\ref{exer:ways to quantify} that are equivalent to others on your list?  
\end{problem}

\begin{exercise}
Suppose that the universe of acceptable values is the set of married people.  Consider the predicate $M(x,y):=``x\mbox{ is married to }y"$.  Discuss the meaning of each of the following.
\begin{enumerate}
\item For all $x$, there exists $y$ such that $M(x,y)$.
\item There exists $y$ such that for all $x$, $M(x,y)$.
\end{enumerate}
\end{exercise}

\begin{exercise}
Suppose that the universe of acceptable values is the set of real numbers.  Consider the predicate $R(x,y):=``x=y^2"$.  Discuss the meaning of each of the following.
\begin{enumerate}
\item There exists $x$ such that there exists $y$ such that $R(x,y)$.
\item There exists $y$ such that there exists $x$ such that $R(x,y)$.
\end{enumerate}
\end{exercise}

\begin{exercise}
Repeat the exercise above but replace the existential quantifiers with universal quantifiers.
\end{exercise}

\begin{problem}
Conjecture a summary of the various possibilities for quantifying predicates involving two variables.  You do \emph{not} need to prove your conjecture.
\end{problem}

\begin{exercise}
Suppose that the universe of acceptable values is the set of real numbers.  Consider the predicate $G(x,y):=``x>y"$.  Find all possible \emph{distinct} ways to bind the variables to create propositions and then determine the truth value of each (you do not need to prove your answers).
\end{exercise}

\end{section}

%%% Local Variables: 
%%% mode: latex
%%% TeX-master: "IntroToProof"
%%% End: 

%\begin{section}{More on Quantification}

In the last section, we introduced the universal quantifier ``for all'' and the existential quantifier ``there exists\ldots such that.''  Here are a couple of important points to remember about quantification:
\begin{enumerate}
\item In order to have a proposition, all variables must be bound.  That is, all variables must be quantified.  This can happen in at least two ways:
\begin{enumerate}
\item The variables are explicitly bound by quantifiers in the same sentence, or
\item The variables are implicitly bound by preceding sentences and/or by context.  \emph{Note:}  Statements of the form ``Let $x=\ldots$" and ``Let $x\in\ldots$" bind the variable $x$ and remove ambiguity.
\end{enumerate}
\item The order of the quantification is important.  Reversing the order of the quantifiers can substantially change the meaning of a proposition.
\end{enumerate}

Using our logical connectives (``and", ``or", ``If\ldots, then\ldots", and ``not") together with quantification, we can build very complex mathematical statements.

\begin{example}\label{ex:def limit}
Let $f$ be a function and consider the formal definition of the limit of $f(x)$ as $x$ approaches $c$:

\begin{quote}
For all $\epsilon >0$, there exists $\delta >0$ such that for all $x$, if $|x-c|<\delta$, then $|f(x)-f(c)|<\epsilon$.
\end{quote}
\end{example}

\begin{exercise}
Identify all the quantifiers from Example \ref{ex:def limit} and any logical connectives.  Are there any implicit bound variables?
\end{exercise}

In order to study the abstract nature of complicated mathematical statements, it is useful to adopt some notation.

\begin{definition}
We use the symbol $\forall$ to denote the universal quantifier ``for all" and the symbol $\exists$ to denote the existential quantifier ``there exists\ldots such that".
\end{definition}

Using our abbreviations for the logical connectives and quantifiers, we can symbolically represent mathematical propositions.

\begin{example}
For each of the following, suppose our universe of discourse is the set of real numbers.
\begin{enumerate}

\item Consider the following (true) proposition:

\begin{quote}
There exists $x$ such that $x^2-1=0$.
\end{quote}

This proposition can be denoted symbolically as $(\exists x)(x^2-1=0)$.

\item Consider the following (false) proposition:

\begin{quote}
For all $x\in \mathbb{N}$, there exists $y\in\mathbb{N}$ such that $y<x$.
\end{quote}

This one can be represented symbolically as $(\forall x)(x\in\mathbb{N}\implies (\exists y)(y\in\mathbb{N}\implies y<x))$ or more simply as $(\forall x\in\mathbb{N})(\exists y\in\mathbb{N})(y<x)$.

\item Consider the following (true) proposition:

\begin{quote}
Every positive real number has a multiplicative inverse.
\end{quote}

There are several ways of representing this statement symbolically.  However, if you unpack what a multiplicative inverse is, you'll get something like $(\forall x)(x>0 \implies (\exists y)(xy=1))$.  Alternatively, you can shorten the statement to $(\forall x>0)(\exists y)(xy=1)$.

\end{enumerate}
\end{example}

\begin{exercise} Convert the following statements into statements using only logical symbols.  Assume that the universe of discourse is the set of real numbers.
\begin{enumerate}
\item There exists a number $x$ such that $x^2+1$ is greater than zero.
\item There exists a natural number $n$ such that $n^2=36$. 
\item For every real number $x$, $x^2$ is greater than or equal to zero.
\end{enumerate}
\end{exercise}

\begin{exercise}
Convert the statement in Example \ref{ex:def limit} into a statement using only logical symbols.
\end{exercise}

\begin{remark}\label{rem:implicit universal}
If $A(x)$ and $B(x)$ are predicates, then it is standard practice for the statement $A(x)\implies B(x)$ to mean $(\forall x)(A(x)\implies B(x))$ (where the universe of discourse for $x$ needs to be made clear).  In this case, we say that the universal quantifier is implicit.
\end{remark}

\begin{exercise}
Find at least two examples earlier in the notes that exhibit the claim made in Remark \ref{rem:implicit universal}.  Attempt to write the statements symbolically.
\end{exercise}

\begin{exercise}
Convert the following proposition into a statement using only logical symbols.  The universe of discourse is the set of real numbers.  (Watch out for implicit quantifiers.)
\begin{quote}
If $\epsilon >0$, then there exists $N\in\mathbb{N}$ such that $1/N<\epsilon$.
\end{quote}
\end{exercise}

\begin{remark}
The symbolic expression $(\forall x)(\forall y)$ can be replaced with the simpler expression $(\forall x,y)$ as long as $x$ and $y$ are coming from the same set.
\end{remark}

\begin{exercise}
For each of the following statements, (i) unpack the statement into words, and (ii) determine whether the statement is true or false.

\begin{enumerate}
\item $(\forall n \in \mathbb{N})(n^2 \geq 5)$
\item $(\exists n \in \mathbb{N})(n^2-1=0)$
\item $(\exists N \in \mathbb{N})(\forall  n > N)(\frac{1}{n} < 0.01)$
\item $(\forall m, n \in \mathbb{Z})(2|m \wedge 2|n \implies 2|(m+n))$
\item $(\forall x \in \mathbb{N})(\exists y \in \mathbb{N})(x-2y=0)$
\item $(\exists x \in \mathbb{N})(\forall y \in \mathbb{N})(y \leq x)$
\end{enumerate}
\end{exercise}

To whet your appetite for the next section, tackle the following questions.

\begin{question}
How would you go about proving a true statement of the form ``For all $x$\ldots"?
\end{question}

\begin{question}
If a statement is false, then its negation is true.  How would you go about negating a statement involving quantifiers?  In particular, what are the negations of $\forall$ and $\exists$, respectively?
\end{question}

\end{section}

%\begin{section}{And Even More on Quantification}

Before we get started, it is important to remind you that we will not be explicitly using the symbolic representation of a given statement in terms of quantifiers and logical connectives.  Nonetheless, having this notation at our disposal allows us to compartmentalize the abstract nature of mathematical propositions and will provide us with a way to talk about the meta-concepts surrounding the construction of proofs.

\begin{definition}
Two quantified propositions are \textbf{equivalent in a given universe of discourse} iff they have the same truth value in that universe.  Two quantified propositions are \textbf{equivalent} iff they are equivalent in every universe of discourse.
\end{definition}

\begin{exercise}
Consider the propositions $(\forall x)(x>3)$ and $(\forall x)(x\geq 4)$.  Are these propositions equivalent if the universe of discourse is the set of integers?  (\emph{Hint:}  What are their truth values in this case?)  Come up with two different universes of discourse that yield different truth values for these propositions.  What can you conclude?
\end{exercise}

\begin{remark}
At this point it is worth pointing out an important distinction.  Consider the propositions ``All cars are red" and ``All natural numbers are positive".  Both of these are instances of the \textbf{logical form} $(\forall x)P(x)$.  It turns out that the first proposition is false and the second is true; however, the logical form is neither true or false.  A logical form is a blueprint for particular propositions.  If we are careful, it makes sense to talk about whether two logical forms are equivalent.  For example, $(\forall x)(P(x)\implies Q(x))$ is equivalent to $(\forall x)(\neg Q(x)\implies \neg P(x))$.  For fixed $P(x)$ and $Q(x)$, these two forms will always have the same truth value independent of the universe of discourse.  If you change $P(x)$ and $Q(x)$, then the truth value may change, but the two forms will still agree.
\end{remark}

\begin{theorem}\label{thm:negation of quantifiers}
Let $P(x)$ be a predicate.  Then
\begin{enumerate}
\item[1.] $\neg (\forall x)P(x)$ is equivalent to $(\exists x)(\neg P(x))$;
\item[2.] $\neg (\exists x)P(x)$ is equivalent to $(\forall x)(\neg P(x))$.
\end{enumerate}
\end{theorem}

\begin{exercise}
Negate each of the following.  Disregard the truth value and the universe of discourse.
\begin{enumerate}
\item $(\forall x)(x>3)$
\item $(\exists x)(x \mbox{ is prime}\wedge x \mbox{ is even})$
\item All cars are red.
\item Every Wookiee is named Chewbacca.
\item Some hippies are republican.
\item For all $x\in\mathbb{N}$, $x^2+x+41$ is prime.
\item There exists $x\in\mathbb{Z}$ such that $1/x\notin\mathbb{Z}$.
\item There does not exist a function $f$ such that if $f$ is continuous, then $f$ is not differentiable.
\end{enumerate}
\end{exercise}

Using Theorem \ref{thm:negation of quantifiers} and our previous results involving quantification, we can negate complex mathematical propositions by working from left to right.

\begin{example}
Consider the proposition
\[
(\exists x\in\mathbb{R})(\forall y\in\mathbb{R})(x+y=0).
\]
It turns out that this statement is false, which means that its negation is true.  That is,
\[
\neg(\exists x\in\mathbb{R})(\forall y\in\mathbb{R})(x+y=0),
\]
which is equivalent to
\[
(\forall x\in\mathbb{R})(\exists y\in\mathbb{R})(x+y\neq 0),
\]
is true.
\end{example}

\begin{example}
Consider the proposition
\[
(\forall x)[x>0\implies (\exists y)(y<0 \wedge xy>0)]
\]
(which happens to be false).  Then
\[
\neg (\forall x)[x>0\implies (\exists y)(y<0 \wedge xy>0)]
\]
is equivalent to
\[
(\exists x)[x>0 \wedge \neg (\exists y)(y<0) \wedge xy>0)],
\]
which is equivalent to 
\[
(\exists x)[x>0 \wedge (\forall y)(y\geq 0 \vee xy\leq 0)].
\]
\end{example}

\begin{exercise}
What previous theorems were used when negating the proposition in the previous example.
\end{exercise}

\begin{exercise}
Negate each of the following.  Disregard the truth value and the universe of discourse.
\begin{enumerate}
\item $(\forall n\in\mathbb{N})(\exists m\in\mathbb{N})(m<n)$
\item $(\forall x,y,z\in\mathbb{Z})((xy \mbox{ is even}\wedge yz\mbox{ is even})\implies xy\mbox{ is even})$
\item For all positive real numbers $x$ there exists a real number $y$ such that $y^2=x$.
\item There exists a married person $x$ such that for all married people $y$, $x$ is married to $y$.
\end{enumerate}
\end{exercise}

At this point, we should be able to use our understanding of quantification to construct counterexamples to complicated false propositions and proofs of complicated true propositions.  Here are some general proof structures for various logical forms.

\begin{skeleton}[Direct Proof of $(\forall x)P(x)$]
Here is what the general structure for a direct proof of the proposition $(\forall x)P(x)$ looks like.

\bigskip

\begin{textbox}
\begin{proof}
Let $x \in U$ (where $U$ is whatever the universe of discourse is).
\begin{center}
$\vdots$\\
(Use definitions and previous results.)\\
$\vdots$
\end{center}
\noindent Therefore, $P(x)$ is true.  Since $x$ was arbitrary, for all $x$, $P(x)$.
\end{proof}
\end{textbox}

\end{skeleton}

\begin{skeleton}[Proof of $(\forall x)P(x)$ by Contradiction]
Here is the general structure for a proof of the proposition $(\forall x)P(x)$ via contradiction.

\bigskip

\begin{textbox}
\begin{proof}
For sake of a contradiction, assume that there exists $x\in U$ (where $U$ is whatever the universe of discourse is) such that $\neg P(x)$.
\begin{center}
$\vdots$\\
(Do something to derive a contradiction.)\\
$\vdots$
\end{center}
\noindent This is a contradiction.  Therefore, for all $x$, $P(x)$ is true.
\end{proof}
\end{textbox}
\end{skeleton}


\begin{skeleton}[Direct Proof of $(\exists x)P(x)$]
Here is what the general structure for a direct proof of the proposition $(\exists x)P(x)$ looks like.

\bigskip

\begin{textbox}
\begin{proof}
(Either use definitions and previous results to deduce that an $x$ exists such that $P(x)$ is true or if you think you have an $x$ that works, just verify that it does.)
\begin{center}
$\vdots$\\
(Do stuff.)\\
$\vdots$
\end{center}Therefore, there exists $x$ such that $P(x)$.
\end{proof}
\end{textbox}

\end{skeleton}

\begin{skeleton}[Proof of $(\exists x)P(x)$ by Contradiction]
Here is the general structure for a proof of the proposition $(\exists x)P(x)$ via contradiction.

\bigskip

\begin{textbox}
\begin{proof}
For sake of a contradiction, assume that for all $x$, $\neg P(x)$.
\begin{center}
$\vdots$\\
(Do something to derive a contradiction.)\\
$\vdots$
\end{center}
\noindent This is a contradiction.  Therefore, there exists $x$ such that $P(x)$.
\end{proof}
\end{textbox}

\end{skeleton}

\begin{question}
Suppose $P(x)$ is a proposition such that $(\forall x)P(x)$ is false.  Which of the above proof situations is identical to providing a counterexample to this proposition?
\end{question}

\begin{remark}
It is important to point out that sometimes we will have to combine various proof techniques in a single proof.  For example, if you wanted to prove a proposition of the form $(\forall x)(P(x) \implies Q(x)$) by contradiction, we would start by assuming that there exists $x$ such that $P(x)$ and $\neg Q(x)$.
\end{remark}

\begin{problem}
For each of the following statements, determine its truth value.  If the statement is false, provide a counterexample.  Prove at least two of the true statements.
\begin{enumerate}
\item For all $n\in\mathbb{N}$, $n^2\geq 5$.

\item There exists $n \in \mathbb{N}$ such that $n^2-1=0$.

%\item For all $x \in \mathbb{N}$, there exists $y \in \mathbb{N}$ such that $x-2y=0$.

\item There exists $x \in \mathbb{N}$ such that for all $y \in \mathbb{N}$, $y \leq x$.

\item For all $x\in\mathbb{Z}$, $x^3\geq x$.

\item For all $n\in\mathbb{Z}$, there exists $m\in\mathbb{Z}$ such that $n+m=0$.

\item There exists integers $a$ and $b$ such that $2a+7b=1$.

\item There do not exist integers $m$ and $n$ such that $2m+4n=7$.

\item For all integers $a, b, c$, if $a$ divides $bc$, then either $a$ divides $b$ or $a$ divides $c$.

%\item For every $m\in\mathbb{Z}$, if $m$ is odd, then there exists $k\in\mathbb{Z}$ such that $m^2=8k+1$.

%\item For all integers $a, b, c,$ and $d$, if $a$ divides $b$ and $a$ divides $c$, then for all integers $x$ and $y$, $a$ divides $bx+cy$.

\end{enumerate}

\end{problem}

\end{section}

\end{chapter}

\begin{comment}
\begin{chapter}{Set Theory and Topology}
%\documentclass[11pt]{article}

\usepackage{amsfonts}
\usepackage{amsmath}
\usepackage{amssymb}
\usepackage{stmaryrd}
\usepackage{amsthm}
\usepackage{fancyhdr}
\usepackage[margin=1in]{geometry}
\usepackage[hang,flushmargin,symbol*]{footmisc}
\usepackage{color}
\definecolor{darkblue}{rgb}{0, 0, .6}
\definecolor{grey}{rgb}{.7, .7, .7}
\usepackage[breaklinks]{hyperref}
\hypersetup{
	colorlinks=true,
	linkcolor=darkblue,
	anchorcolor=darkblue,
	citecolor=darkblue,
	pagecolor=darkblue,
	urlcolor=darkblue,
	pdftitle={},
	pdfauthor={}
}

\pagestyle{fancy}

\lhead{\scriptsize Notes for an Introduction to Proof Course (Version Spring 2013)} 
\rhead{\scriptsize Instructor: \href{http://danaernst.com}{D.C. Ernst}}
\lfoot{\scriptsize This work is an adaptation of notes written by Stan Yoshinobu of Cal Poly and Matthew Jones of California State University, Dominguez Hills.} 
\cfoot{}
\renewcommand{\headrulewidth}{0.4pt} 
\renewcommand{\footrulewidth}{0.4pt} 

\theoremstyle{definition}
\newtheorem{theorem}{Theorem}[section]
\newtheorem{acknowledgement}[theorem]{Acknowledgement}
\newtheorem{algorithm}[theorem]{Algorithm}
\newtheorem{axiom}[theorem]{Axiom}
\newtheorem{case}[theorem]{Case}
\newtheorem{claim}[theorem]{Claim}
\newtheorem{conclusion}[theorem]{Conclusion}
\newtheorem{condition}[theorem]{Condition}
\newtheorem{conjecture}[theorem]{Conjecture}
\newtheorem{corollary}[theorem]{Corollary}
\newtheorem{criterion}[theorem]{Criterion}
\newtheorem{definition}[theorem]{Definition}
\newtheorem{example}[theorem]{Example}
\newtheorem{exercise}[theorem]{Exercise}
\newtheorem{journal}[theorem]{Journal}
\newtheorem{lemma}[theorem]{Lemma}
\newtheorem{notation}[theorem]{Notation}
\newtheorem{problem}[theorem]{Problem}
\newtheorem{proposition}[theorem]{Proposition}
\newtheorem{remark}[theorem]{Remark}
\newtheorem{solution}[theorem]{Solution}
\newtheorem{summary}[theorem]{Summary}
\newtheorem{question}[theorem]{Question}

\begin{document}

\addtocounter{section}{1}

\begin{section}{Introduction to Set Theory and Topology}

At its essence, all of mathematics is built on set theory.  In this chapter, we will introduce some of the basics of sets and their properties.

\begin{subsection}{Sets} 

\begin{definition}
A \textbf{set} is a collection of objects called \textbf{elements}. If $A$ is a set and $x$ is an element of $A$, we write
$x\in A$. Otherwise, we write $x\notin A$.
\end{definition}

\begin{definition}
The set containing no elements is called the \textbf{empty set}, and is denoted by the symbol $\emptyset$.
\end{definition}

If we think of a set as a box containing some stuff, then the empty set is a box with nothing in it.

\begin{definition} 
The language associated to sets is specific.  We will often define sets using the following notation, called \textbf{set builder notation}.
\[
S=\{x \in A: x \mbox{ satisfies some condition}\}
\]
The first part ``$x \in A$" denotes what type of $x$ is being considered.  The statements to the right of the colon are the conditions that $x$ must satisfy in order to be members of the set.  This notation is read as ``The set of all $x$ in $A$ such that $x$ satisfies some condition,'' where ``some condition" is something specific about the restrictions on $x$ relative to $A$.
\end{definition}

\begin{exercise}
Unpack each of the following sets into a description using a sentence and see if you can determine exactly what elements each set contains.
\begin{enumerate}
\item $M=\{x \in \mathbb{R} :  x \geq 2 \}$
\item $A=\{x \in \mathbb{N} : x = 3k \mbox{ for some } k\in \mathbb{N} \}$
\item $T=\{t \in \mathbb{R} : t^2 \leq 2 \}$
\item $H=\{t \in \mathbb{R} : t = 1 - \frac{1}{n} \mbox{, where } n \in \mathbb{Z} \}$
\end{enumerate}
\end{exercise}

\begin{exercise}
Write each of the following sentences using set builder notation.
\begin{enumerate}
\item Suppose $R$ is the set of all real numbers $x$ such that $x$ is less than $-\sqrt{2}$. 
\item Suppose $A$ is the set of all real numbers $y$, such that $y$ is greater than $-12$ and less than 42.4.
\item Suppose $D$ is the set of all even natural numbers.
\end{enumerate}
\end{exercise}

\begin{definition}
If $A$ and $B$ are sets, then we say that $A$ is a \textbf{subset} of $B$, written $A\subseteq B$, provided that every element of $A$ is also an element of $B$.
\end{definition}

\begin{remark}
Observe that $A\subseteq B$ is equivalent to ``For all $x$ (in the universe of discourse), if $x\in A$, then $x\in B$."  Since we know how to deal with ``for all" statements and conditional propositions, we know how to go about proving $A\subseteq B$.
\end{remark}

\begin{question}
Suppose that $A$ and $B$ are sets.  Describe a general strategy for proving that $A\subseteq B$.
\end{question}

\begin{theorem}
Let $S$ be a set.  Then
\begin{enumerate}
\item $S\subseteq S$,
\item $\emptyset \subseteq S$.
\end{enumerate}
\end{theorem}

\begin{exercise}
List all of the subsets of $A=\{1,2,3\}$.  Any conjectures about how many there might be for a set with $n$ elements?
\end{exercise}

\begin{theorem}[*, Transitivity of subsets]
Suppose that $A$, $B$, and $C$ are sets.  If $A\subseteq B$ and $B\subseteq C$, then $A\subseteq C$.
\end{theorem}

\begin{definition}
If $A\subseteq B$, then $A$ is called a \textbf{proper subset} provided that $A\neq B$.  In this case, we may write $A\subset B$ or $A\subsetneq B$.\footnote{\emph{Warning:} Some books use $\subset$ to mean $\subseteq$.}
\end{definition}

\begin{definition}[Interval Notation]
For $a,b\in\mathbb{R}$ with $a<b$, we define the following.
\begin{enumerate}
\item $(a,b)=\{x\in\mathbb{R}:a<x<b\}$
\item $(a,\infty)=\{x\in\mathbb{R}:a<x\}$
\item $(-\infty,b)=\{x\in\mathbb{R}:x<b\}$
\item $[a,b]=\{x\in\mathbb{R}:a\leq x\leq b\}$
\end{enumerate}
We analogously define $[a,b)$, $(a,b]$, $[a,\infty)$, and $(-\infty,b]$.
\end{definition}

\begin{exercise}
Provide two examples of proper subsets of the interval $[0,1]$.
\end{exercise}

Here are some more definitions.  In each case, take $U$ to be the universe of discourse.

\begin{definition}
The \textbf{union} of the sets $A$ and $B$ is $A \cup B =\{x\in U : x\in A \mbox{ or } x\in B \}$.
\end{definition}

\begin{definition}
The \textbf{intersection} of the sets $A$ and $B$ is $A \cap B =\{x\in U : x\in A \mbox{ and } x\in B \}$.
\end{definition}

\begin{definition}
The \textbf{set difference} of the sets $A$ and $B$ is $A \setminus B =\{x\in U : x\in A \mbox{ and } x\notin B \}$.
\end{definition}

\begin{definition}
The \textbf{complement of $A$} (relative to $U$) is the set $A^c=U \setminus A =\{x \in U : x \notin A\}$.
\end{definition}

\begin{definition}
If two sets $A$ and $B$ have the property that $A \cap B = \emptyset$, then we say that $A$ and $B$ are \textbf{disjoint} sets.
\end{definition}

\begin{exercise}
Suppose that the universe of discourse is $U=\{1,2,3,4,5,6,7,8,9,10\}$.  Let $A=\{1, 2, 3, 4, 5\}$, $B=\{1, 3, 5\}$, and $C=\{2, 4, 6, 8\}$.  Find each of the following.
\begin{enumerate}
  \item $A \cap C$
  \item $A \cap B$
  \item $A \cup C$
  \item $A \cup B$
  \item $A\setminus B$
  \item $B \setminus A$
  \item $C \setminus B$
  \item $B \cap C$
  \item $B^c$
  \item $A^c$
  \item $(A\cup B)^c$
  \item $A^c\cap B^c$
\end{enumerate}
\end{exercise}

\begin{exercise}
Suppose that the universe of discourse is $U=\mathbb{R}$.  Let $A=[-3,-1)$, $B=(-2.5,2)$, and $C=(-2,0]$.  Find each of the following.
\begin{enumerate}
\item $A^c$
  \item $A \cap C$
  \item $A \cap B$
  \item $A \cup C$
  \item $A \cup B$
  \item $(A\cap B)^c$
  \item $(A\cup B)^c$
  \item $A \setminus B$
  \item $A\setminus (B \cup C)$
  \item $B \setminus A$
  \item $B \cap C$
\end{enumerate}
\end{exercise}

\begin{theorem}[*]
Let $A$ and $B$ be sets.  If $A \subseteq B$, then $B^c \subseteq A^c$.
\end{theorem}

\begin{definition}
Two sets $A$ and $B$ are \textbf{equal} if and only if $A \subseteq B$ and $B \subseteq A$.  In this case we write $A = B$.
\end{definition}

\begin{remark}
Given two sets $A$ and $B$, if we want to prove $A=B$, then we have to do two separate ``mini" proofs: one for $A\subseteq B$ and one for $B\subseteq A$.
\end{remark}

\begin{theorem}[*]
Let $A$ and $B$ be sets.  Then $A\setminus B = A \cap B^c$.
\end{theorem}

\begin{theorem}[*, DeMorgan's Law]
Let $A$ and $B$ be sets. Then 
\begin{enumerate}
\item $(A \cup B)^c = A^c \cap B^c$,
\item $(A \cap B)^c = A^c \cup B^c$.
\end{enumerate}
(You only need to prove one of these; the other is similar.)
\end{theorem}

\begin{theorem}[*, Distribution of Union and Intersection]
Let $A$, $B$, and $C$ be sets. Then
\begin{enumerate}
\item $A \cup(B\cap C) = (A\cup B)\cap (A\cup C)$,
\item $A\cap (B\cup C)= (A\cap B)\cup (A\cap C)$.
\end{enumerate}
(You only need to prove one of these; the other is similar.)
\end{theorem}

\end{subsection}

\end{section}

\end{document}
%\documentclass[11pt]{article}

\usepackage{amsfonts}
\usepackage{amsmath}
\usepackage{amssymb}
\usepackage{stmaryrd}
\usepackage{amsthm}
\usepackage{fancyhdr}
\usepackage[margin=1in]{geometry}
\usepackage[hang,flushmargin,symbol*]{footmisc}
\usepackage{color}
\definecolor{darkblue}{rgb}{0, 0, .6}
\definecolor{grey}{rgb}{.7, .7, .7}
\usepackage[breaklinks]{hyperref}
\hypersetup{
	colorlinks=true,
	linkcolor=darkblue,
	anchorcolor=darkblue,
	citecolor=darkblue,
	pagecolor=darkblue,
	urlcolor=darkblue,
	pdftitle={},
	pdfauthor={}
}

\pagestyle{fancy}

\lhead{\scriptsize Course Notes for Introduction to Proof (Version 1.1)} 
\rhead{\scriptsize Instructor: \href{http://danaernst.com}{D.C. Ernst}}
\lfoot{\scriptsize This work is an adaptation of notes written by Stan Yoshinobu of Cal Poly and Matthew Jones of California State University, Dominguez Hills.} 
\cfoot{}
\renewcommand{\headrulewidth}{0.4pt} 
\renewcommand{\footrulewidth}{0.4pt} 

\theoremstyle{definition}
\newtheorem{theorem}{Theorem}[section]
\newtheorem{acknowledgement}[theorem]{Acknowledgement}
\newtheorem{algorithm}[theorem]{Algorithm}
\newtheorem{axiom}[theorem]{Axiom}
\newtheorem{case}[theorem]{Case}
\newtheorem{claim}[theorem]{Claim}
\newtheorem{conclusion}[theorem]{Conclusion}
\newtheorem{condition}[theorem]{Condition}
\newtheorem{conjecture}[theorem]{Conjecture}
\newtheorem{corollary}[theorem]{Corollary}
\newtheorem{criterion}[theorem]{Criterion}
\newtheorem{definition}[theorem]{Definition}
\newtheorem{example}[theorem]{Example}
\newtheorem{exercise}[theorem]{Exercise}
\newtheorem{journal}[theorem]{Journal}
\newtheorem{lemma}[theorem]{Lemma}
\newtheorem{notation}[theorem]{Notation}
\newtheorem{problem}[theorem]{Problem}
\newtheorem{proposition}[theorem]{Proposition}
\newtheorem{remark}[theorem]{Remark}
\newtheorem{solution}[theorem]{Solution}
\newtheorem{summary}[theorem]{Summary}
\newtheorem{question}[theorem]{Question}

\begin{document}

\addtocounter{section}{1}

\begin{section}{Set Theory and Topology (Continued)}

\addtocounter{subsection}{1}
\addtocounter{theorem}{27}

\begin{subsection}{Power Sets}

We've already seen that using union, intersection, set difference, and complement that we can create new sets (in the same universe) from existing sets.  In this section, we will describe another way to generate new sets; however, the new sets will not ``live" in the same universe this time.

\begin{definition}
If $S$ is a set, then the \textbf{power set} of $S$ is the set of subsets of $S$.  The power set of $S$ is denoted $\mathcal{P}(S)$.
\end{definition}

\begin{remark}
It follows immediately from the definition that $A\subseteq S$ iff $A\in\mathcal{P}(S)$.\footnote{Recall that ``iff" is an abbreviation for `if and only if", which is a statement of the form $A\iff B$ for propositions $A$ and $B$.  Recall that this is short for both $A\implies B$ \emph{and} $B\implies A$.}
\end{remark}

\begin{example}
If $S=\{a,b\}$, then $\mathcal{P}=\{\emptyset, \{a\}, \{b\}, S\}$.
\end{example}

\begin{question}
Implicit in the definition of power set is that $S$ is a subset of some fixed universe $U$.  In what universe does $\mathcal{P}(S)$ live?
\end{question}

\begin{exercise}
For each of the following sets, find the power set.
\begin{enumerate}
\item $W=\{\circ, \triangle, \square\}$
\item $O=\{a,\{a\}\}$
\item $R=\emptyset$
\item $D=\{\emptyset\}$
\end{enumerate}
\end{exercise}

\begin{conjecture}
How many subsets do you think that a set with $n$ elements has?  What if $n=0$?  You do not need to prove your conjecture at this time.  We will prove this later using mathematical induction.
\end{conjecture}

\begin{exercise}
Do your best to describe $\mathcal{P}(\mathbb{N})$.  You cannot write down all of $\mathcal{P}(\mathbb{N})$.  Why not?
\end{exercise}

\begin{remark}
It is important to realize that the concepts of \emph{element} and \emph{subset} need to be carefully delineated.  For example, consider the set $A=\{x,y\}$.  The object $x$ is an element of $A$, but the object $\{x\}$ is both a subset of $A$ and an element of $\mathcal{P}(A)$.  This can get confusing rather quickly.  Consider the set $O$ from the previous example.  The set $\{a\}$ happens to be an element of $O$, a subset of $O$, and an element of  $\mathcal{P}(O)$.
\end{remark}

\begin{theorem}[*]
Let $S$ and $T$ be sets.  Then $S\subseteq T$ iff $\mathcal{P}(S)\subseteq \mathcal{P}(T)$.\footnote{To prove this theorem, you have to write two distinct subproofs: $A\implies B$ and $B\implies A$.}
\end{theorem}

\begin{theorem}[*]
Let $S$ and $T$ be sets.  Then $\mathcal{P}(S)\cap\mathcal{P}(T)=\mathcal{P}(S\cap T)$.
\end{theorem}

\begin{theorem}[*]
Let $S$ and $T$ be sets.  Then $\mathcal{P}(S)\cup\mathcal{P}(T)\subseteq \mathcal{P}(S\cup T)$.
\end{theorem}


\begin{exercise}
Let $S$ and $T$ be sets.
\begin{enumerate}
\item Provide a counterexample to show that it is not necessarily true that $\mathcal{P}(S)\cup\mathcal{P}(T)= \mathcal{P}(S\cup T)$.
\item Is it ever true that $\mathcal{P}(S)\cup\mathcal{P}(T)= \mathcal{P}(S\cup T)$ or are $\mathcal{P}(S)\cup\mathcal{P}(T)$ and $\mathcal{P}(S\cup T)$ always different sets?
\end{enumerate}
\end{exercise}

We now turn out attention to the issue of whether there is one mother of all universal sets.  Before reading any further, consider this for a moment.  That is, is there one largest set that all other sets are a subset of?  Or, in other words, is there a set of all sets?  To help wrap our heads around this issue, consider the following riddle.

\begin{quote}
In Seville, the barber shaves all those and only those who do not shave themselves.  Who shaves the barber?
\end{quote}

\begin{problem}
Discuss the ``Barber of Seville" riddle.  Does the barber shave himself or not?
\end{problem}

The ``Barber of Seville" riddle is an example of a \textbf{paradox}.  Now, suppose that there is a set of all sets and call it $\mathcal{U}$.  Then we can write $\mathcal{U}=\{A:A\mbox{ is a set}\}$.

\begin{problem}
Given our definition of $\mathcal{U}$, explain why it is an element of itself.
\end{problem}

If we continue with this line of reasoning, it must be the case that some sets are elements of themselves and some are not.  Let $X$ be the set of all sets that are elements of themselves and let $Y$ be the set of all sets that are not elements of themselves.

\begin{question}
Does $Y$ belong to $X$ or $Y$?  Explain why this is a paradox.
\end{question}

The above paradox is one way of phrasing a paradox referred to as \textbf{Russell's paradox}.  Okay, how did we get into this mess in the first place?!  By assuming the existence of a set of all sets, we can produce all sorts of paradoxes.  The only way to avoid the paradoxes is to conclude that there is no set of all sets.  Here is some more evidence that we shouldn't assume the existence of a set of all sets.

\begin{question}
If $\mathcal{U}$ is the set of all sets, then what is the relationship between $\mathcal{U}$ and $\mathcal{P}(\mathcal{U})$?  What about $\mathcal{P}(\mathcal{P}(\mathcal{U})$?
\end{question}

The upshot is that the collection of all sets is \emph{not} a set!  Moreover, there are lots of collections that are not sets.

\end{subsection}

\end{section}

\end{document}
%\documentclass[11pt]{article}

\usepackage{amsfonts}
\usepackage{amsmath}
\usepackage{amssymb}
\usepackage{stmaryrd}
\usepackage{amsthm}
\usepackage{fancyhdr}
\usepackage[margin=1in]{geometry}
\usepackage[hang,flushmargin,symbol*]{footmisc}
\usepackage{color}
\definecolor{darkblue}{rgb}{0, 0, .6}
\definecolor{grey}{rgb}{.7, .7, .7}
\usepackage[breaklinks]{hyperref}
\hypersetup{
	colorlinks=true,
	linkcolor=darkblue,
	anchorcolor=darkblue,
	citecolor=darkblue,
	pagecolor=darkblue,
	urlcolor=darkblue,
	pdftitle={},
	pdfauthor={}
}

\pagestyle{fancy}

\lhead{\scriptsize Course Notes for Introduction to Proof (Version 1.1)} 
\rhead{\scriptsize Instructor: \href{http://danaernst.com}{D.C. Ernst}}
\lfoot{\scriptsize This work is an adaptation of notes written by Stan Yoshinobu of Cal Poly and Matthew Jones of California State University, Dominguez Hills.} 
\cfoot{}
\renewcommand{\headrulewidth}{0.4pt} 
\renewcommand{\footrulewidth}{0.4pt} 

\theoremstyle{definition}
\newtheorem{theorem}{Theorem}[section]
\newtheorem{acknowledgement}[theorem]{Acknowledgement}
\newtheorem{algorithm}[theorem]{Algorithm}
\newtheorem{axiom}[theorem]{Axiom}
\newtheorem{case}[theorem]{Case}
\newtheorem{claim}[theorem]{Claim}
\newtheorem{conclusion}[theorem]{Conclusion}
\newtheorem{condition}[theorem]{Condition}
\newtheorem{conjecture}[theorem]{Conjecture}
\newtheorem{corollary}[theorem]{Corollary}
\newtheorem{criterion}[theorem]{Criterion}
\newtheorem{definition}[theorem]{Definition}
\newtheorem{example}[theorem]{Example}
\newtheorem{exercise}[theorem]{Exercise}
\newtheorem{journal}[theorem]{Journal}
\newtheorem{lemma}[theorem]{Lemma}
\newtheorem{notation}[theorem]{Notation}
\newtheorem{problem}[theorem]{Problem}
\newtheorem{proposition}[theorem]{Proposition}
\newtheorem{remark}[theorem]{Remark}
\newtheorem{solution}[theorem]{Solution}
\newtheorem{summary}[theorem]{Summary}
\newtheorem{question}[theorem]{Question}

\begin{document}

\addtocounter{section}{1}

\begin{section}{Set Theory and Topology (Continued)}

\addtocounter{subsection}{2}
\addtocounter{theorem}{43}

\begin{subsection}{Indexing Sets}

Suppose we wish to consider the following collection of open intervals:
\[
(0,1), (0,1/2), (0,1/4), \ldots, (0,1/2^{n-1}), \ldots
\]
This collection has a natural way for us to ``index" the sets:
\[
I_1=(0,1), I_2=(0,1/2), \ldots, I_n=(0,1/2^{n-1}), \ldots
\]
In this case the sets are \textbf{indexed} by the set $\mathbb{N}$.  The subscripts on the capital letters are taken from the \textbf{index set}.  If we wanted to talk about an arbitrary set from this indexed collection, we could use the notation $I_n$.

Let's consider another example:
\[
\{a\}, \{a,b\}, \{a,b,c\}, \ldots, \{a,b,c,\ldots,z\}
\]
An obvious way to index these sets is as follows:
\[
A_1=\{a\}, A_2=\{a,b\}, A_3=\{a,b,c\}, \ldots, A_{26}=\{a,b,c,\ldots,z\}
\]
In this case, the collection of sets is indexed by $\{1,2,\ldots, 26\}$.

\begin{remark}
Using indexing sets in mathematics is an extremely useful notational tool, but it is important to keep straight the difference between the sets that are being indexed, the elements in each set being indexed, the indexing set, and the elements of the indexing set.
\end{remark}

Any set (finite or infinite) can be used as an indexing set.  Often capital Greek letters are used to denote arbitrary indexing sets and small Greek letters to represent elements of these sets.  For example, we might use $\Delta$ (capital delta) to refer to an indexing set and write $\alpha \in \Delta$ for an individual index.  Typically, if the indexing set is some subset of $\mathbb{Z}$ (like $\mathbb{N}$), then we would use letters like $k,m,n,l$ for an individual index.  Likewise, if the indexing set is $\mathbb{R}$, then we might use $s,t,x,y$ as indices.  

\begin{example}
Here are some examples of common notation that you will encounter.
\begin{enumerate}
\item If $\Delta$ is a set and we have a collection of sets indexed by $\Delta$, then we may write
\[
\{S_{\alpha}\}_{\alpha\in \Delta}
\]
to refer to this collection.  We read this as ``the set of $S$-alphas over alpha in Delta."
\item If a collection of sets is indexed by the natural numbers, then we may write
\[
\{U_n\}_{n\in\mathbb{N}}
\]
or
\[
\{U_n\}_{n=1}^{\infty}.
\]
\item Borrowing from this idea, we can write the collection $\{A_1,\ldots,A_{26}\}$ from the beginning of the section as
\[
\{A_n\}_{n=1}^{26}.
\]

\end{enumerate}
\end{example}

\begin{definition}
Suppose we have a collection $\{A_{\alpha}\}_{\alpha\in\Delta}$.

\begin{enumerate}
\item The \textbf{union of the entire collection} is defined via
\[
\bigcup_{\alpha\in\Delta} A_{\alpha}=\{x:x\in A_{\alpha} \mbox{ for some }\alpha\in \Delta\}.
\]

\item The \textbf{intersection of the entire collection} is defined via
\[
\bigcap_{\alpha\in\Delta} A_{\alpha}=\{x:x\in A_{\alpha} \mbox{ for all }\alpha\in \Delta\}.
\]

\end{enumerate}
\end{definition}

\begin{remark} 
In the special case that $\Delta=\mathbb{N}$, we write
\[
\bigcup_{n=1}^{\infty}A_n= \{ x : x \in A_n \mbox{ for some } n \in \mathbb{N}\}
\] 
and
\[
\bigcap_{n=1}^{\infty}A_n= \{ x : x \in A_n \mbox{ for all } n \in \mathbb{N}\}.
\] 
\end{remark}

\begin{exercise}
Let $\{I_n\}_{n\in\mathbb{N}\cup\{0\}}$ be the collection of open intervals from the beginning of the section.  Find each of the following.
\begin{enumerate}
\item $\displaystyle \bigcup_{n\in\mathbb{N}\cup\{0\}}I_n$
\item $\displaystyle \bigcap_{n\in\mathbb{N}\cup\{0\}}I_n$
\end{enumerate}
\end{exercise}

\begin{exercise}
Repeat the previous exercise, but assume that the sets are closed intervals.
\end{exercise}

\begin{exercise}
Let $\{A_n\}_{n=1}^{26}$ be the collection from earlier in the section.  Find each of the following.
\begin{enumerate}
\item $\displaystyle \bigcup_{n=1}^{26}A_n$
\item $\displaystyle \bigcap_{n=1}^{26}A_n$
\end{enumerate}
\end{exercise}

\begin{exercise}
Let $S_n = \{x \in \mathbb{R} \ : \ n-1<x<n \}$.  Find each of the following.
\begin{enumerate}
\item $\displaystyle \bigcup_{n=1}^{\infty}S_n$

\item $\displaystyle \bigcap_{n=1}^{\infty}S_n$
\end{enumerate}
\end{exercise}

\begin{exercise}
Let $T_n = \{x \in \mathbb{R} \ : \ -\frac{1}{n}<x< \frac{1}{n} \}$.  Find each of the following.
\begin{enumerate}
\item $\displaystyle \bigcup_{n=1}^{\infty}T_n$

\item $\displaystyle \bigcap_{n=1}^{\infty}T_n$
\end{enumerate}

\end{exercise}

\begin{exercise}
For each $r\in\mathbb{Q}$ (the rational numbers), let $N_r$ be the set containing all real numbers \emph{except} $r$.  Find each of the following.
\begin{enumerate}
\item $\displaystyle \bigcup_{r\in\mathbb{Q}}N_r$

\item $\displaystyle \bigcap_{r\in\mathbb{Q}}N_r$
\end{enumerate}

\end{exercise}

\begin{definition}
We say that a collection of sets $\{A_{\alpha}\}_{\alpha\in\Delta}$ is \textbf{pairwise disjoint} if $A_{\alpha} \cap A_{\beta}=\emptyset$ whenever $\alpha\neq \beta$.
\end{definition}

\begin{exercise}
Draw a Venn diagram of a collection of 3 sets that are pairwise disjoint.
\end{exercise}

\begin{exercise}
Provide an example of a collection of three sets, say $\{A_1, A_2, A_3\}$, such that the collection is \emph{not} pairwise disjoint, but 
\[
\bigcap_{n=1}^3 A_n=\emptyset.
\]
\end{exercise}

\begin{theorem}[*, Generalized Distribution of Union and Intersection]
Suppose we have a collection $\{A_{\alpha}\}_{\alpha\in\Delta}$.  Let $B$ be any set.  Then
\begin{enumerate}
\item $\displaystyle B \cup \left(\bigcap_{\alpha\in\Delta}A_{\alpha}\right)=\bigcap_{\alpha\in\Delta}(B\cup A_{\alpha})$,
\item $\displaystyle B \cap \left(\bigcup_{\alpha\in\Delta}A_{\alpha}\right)=\bigcup_{\alpha\in\Delta}(B\cap A_{\alpha})$.
\end{enumerate}
(You only need to prove one of these; the other is similar.)
\end{theorem}

\begin{theorem}[*, Generalized DeMorgan's Law]
Suppose we have a collection $\{A_{\alpha}\}_{\alpha\in\Delta}$.  Then
\begin{enumerate}
\item $\displaystyle \left(\bigcup_{\alpha\in\Delta} A_{\alpha}\right)^C=\bigcap_{\alpha\in\Delta}A_{\alpha}^{C}$,
\item $\displaystyle \left(\bigcap_{\alpha\in\Delta} A_{\alpha}\right)^C=\bigcup_{\alpha\in\Delta}A_{\alpha}^{C}$.
\end{enumerate}
(You only need to prove one of these; the other is similar.)
\end{theorem}

\end{subsection}

\end{section}

\end{document}
%\documentclass[11pt]{article}

\usepackage{amsfonts}
\usepackage{amsmath}
\usepackage{amssymb}
\usepackage{stmaryrd}
\usepackage{amsthm}
\usepackage{fancyhdr}
\usepackage[margin=1in]{geometry}
\usepackage[hang,flushmargin,symbol*]{footmisc}
\usepackage{color}
\definecolor{darkblue}{rgb}{0, 0, .6}
\definecolor{grey}{rgb}{.7, .7, .7}
\usepackage[breaklinks]{hyperref}
\hypersetup{
	colorlinks=true,
	linkcolor=darkblue,
	anchorcolor=darkblue,
	citecolor=darkblue,
	pagecolor=darkblue,
	urlcolor=darkblue,
	pdftitle={},
	pdfauthor={}
}

\pagestyle{fancy}

\lhead{\scriptsize Course Notes for Logic, Proof, \& Axiomatic Systems (Spring 2011)} 
\rhead{\scriptsize Instructor: \href{http://oz.plymouth.edu/~dcernst}{D.C. Ernst}} 
\lfoot{\scriptsize This work is an adaptation of notes written by Stan Yoshinobu of Cal Poly and Matthew Jones of California State University, Dominguez Hills.} 
\cfoot{} 
\renewcommand{\headrulewidth}{0.4pt} 
\renewcommand{\footrulewidth}{0.4pt} 

\theoremstyle{definition}
\newtheorem{theorem}{Theorem}[section]
\newtheorem{acknowledgement}[theorem]{Acknowledgement}
\newtheorem{algorithm}[theorem]{Algorithm}
\newtheorem{axiom}[theorem]{Axiom}
\newtheorem{case}[theorem]{Case}
\newtheorem{claim}[theorem]{Claim}
\newtheorem{conclusion}[theorem]{Conclusion}
\newtheorem{condition}[theorem]{Condition}
\newtheorem{conjecture}[theorem]{Conjecture}
\newtheorem{corollary}[theorem]{Corollary}
\newtheorem{criterion}[theorem]{Criterion}
\newtheorem{definition}[theorem]{Definition}
\newtheorem{example}[theorem]{Example}
\newtheorem{exercise}[theorem]{Exercise}
\newtheorem{journal}[theorem]{Journal}
\newtheorem{lemma}[theorem]{Lemma}
\newtheorem{notation}[theorem]{Notation}
\newtheorem{problem}[theorem]{Problem}
\newtheorem{proposition}[theorem]{Proposition}
\newtheorem{remark}[theorem]{Remark}
\newtheorem{solution}[theorem]{Solution}
\newtheorem{summary}[theorem]{Summary}
\newtheorem{question}[theorem]{Question}

\begin{document}

\addtocounter{section}{1}

\begin{section}{Set Theory and Topology}

\addtocounter{subsection}{3}
\addtocounter{theorem}{58}

\begin{subsection}{Basic Topology of $\mathbb{R}$}

\begin{remark}
For this entire section, our universe of discourse is the set of real numbers.  You may assume all the usual basic algebraic properties of the real numbers (addition, subtraction, multiplication, division, commutative property, distribution, etc.).
\end{remark}

Recall that an \textbf{axiom} is a statement that we \emph{assume} to be true.  Here are some useful axioms of the real numbers.

\begin{axiom} If $p$ and $q$ are two different real numbers in $\mathbb{R}$, then there is a number between them. %(For example $(p+q)/2$.) 
\end{axiom}

\begin{exercise}
Given real numbers $p$ and $q$ with $p<q$, construct a real number $x$ such that $p<x<q$.  (We know such a point must exist by the previous example, but this exercise is asking you to produce an actual candidate.)
\end{exercise}

\begin{axiom} (Linear ordering) If $a$, $b$, and $c$ are real numbers, then:
\begin{enumerate}
\item If $a < b$ and $b<c$, then $a<c$;
\item Exactly one of the following is true: (i) $a < b$, (ii) $a=b$, or (iii) $a>b$.
\end{enumerate}
\end{axiom}

\begin{axiom} If $p$ is a real number, then there exists $q,r\in\mathbb{R}$ such that $q<p<r$.
\end{axiom}

\begin{axiom} (Archimedean Property) If $x$ is a real number, then either (i) $x$ is an integer or (ii) there exists an integer $n$, such that $n<x<n+1$. 
\end{axiom}

\begin{definition}
Suppose $a,b\in\mathbb{R}$ such that $a<b$.  The intervals $(a,b), (-\infty,b), (a,\infty)$ are called \textbf{open intervals} while the interval $[a,b]$ is called a \textbf{closed interval}.
\end{definition}

\begin{remark} 
In this class, we will always assume that any time we write $(a,b), [a,b], (a,b]$, or $[a,b)$ that $a<b$. 
\end{remark}

\begin{exercise} Give an example of each of the following:
\begin{enumerate}
\item an open interval
\item a closed interval
\item an interval that is neither open or closed
\item an infinite set that is not an interval
\end{enumerate}
\end{exercise}

\begin{definition} A set $U$ is called an \textbf{open set} iff for every $t \in U$, there exists an open interval containing $t$ such that the open interval is a subset of $U$.  We define the empty set to be open.\end{definition}

%where should this go?
%\begin{proposition} 
%Given a number $\varepsilon >0$, there exists a natural number $N$ such that $\frac{1}{N}<\varepsilon$.
%\end{proposition}

\begin{problem} 
Prove that the set $I=(1,2)$ is an open set.
\end{problem}

\begin{theorem}[*]
Every open interval is an open set. 
\end{theorem}

\begin{theorem}
The real numbers form an open set.
\end{theorem}

\begin{theorem}[*]
Every closed interval is not an open set.
\end{theorem}

\begin{theorem}
Let $x\in\mathbb{R}$.  Then the set $\{x\}$ is not open.
\end{theorem}

\begin{exercise} 
Determine whether $\{4,17,42\}$ is an open set, and briefly justify your assertion. 
\end{exercise}

\begin{theorem}[*]\label{finite union of open sets}
Let $A$ and $B$ be open sets.  Then $A\cup B$ is an open set. 
\end{theorem}

\begin{theorem}[*]\label{finite intersection of open sets}
Let $A$ and $B$ be open sets.  Then $A\cap B$ is an open set.
\end{theorem}

\begin{theorem}[*]\label{union of open sets}
Let $\{U_{\alpha}\}_{\alpha\in\Delta}$ be a collection of open sets.  Then
\[
\bigcup_{\alpha\in\Delta} U_{\alpha}
\]
is an open set.
\end{theorem}

\begin{exercise}\label{intersection of open sets}\
\begin{enumerate}
\item Find a collection of open sets $\{U_{\alpha}\}_{\alpha\in\Delta}$ such that
\[
\bigcap_{\alpha\in\Delta} U_{\alpha}
\]
is not an open set.
\item Find a collection of open sets $\{B_{\alpha}\}_{\alpha\in\Delta}$ such that
\[
\bigcap_{\alpha\in\Delta} B_{\alpha}
\]
is an open set.
\end{enumerate}
\end{exercise}

\begin{remark}
Taken together, Theorems \ref{finite union of open sets}--\ref{union of open sets} and Exercise \ref{intersection of open sets} tell us that the union of open sets is open, but that the intersection of open sets may or may not be open.  However, if we are taking the intersection of finitely many open sets, then the intersection will be open.
\end{remark}

\begin{exercise}
Provide an example of an open set that is not an open interval.
\end{exercise}

\begin{exercise}
Determine whether each of the following sets are open or not open.
\begin{enumerate}
\item $\displaystyle W=\bigcup_{n=2}^{\infty} \left(n - \frac{1}{2},n\right)$
\item $\displaystyle X=\bigcap_{n=1}^{\infty} \left(-\frac{1}{n}, \frac{1}{n}\right)$
\end{enumerate}
\end{exercise}

\begin{definition}
A point $p$ is a \textbf{limit point of the set $S$} iff for every open interval $I$ containing $p$, there exists a point $q \in I$ such that $q \in S$ with $q\neq p$.
\end{definition}

\begin{problem}
Consider the open interval $S=(1,2)$. Prove each of the following.
\begin{enumerate}
\item The point $p=2$ is a limit point of $S$.
\item If $p\in S$, then $p$ is a limit point of $S$.
\item If $p<1$ or $p>2$, then $p$ is not a limit point of $S$.
\end{enumerate}
\end{problem}

\begin{theorem}[*]
A point $p$ is a limit point of $(a,b)$ iff $p\in [a,b]$.
\end{theorem}

\begin{problem}
Prove that the point $p=0$ is a limit point of $S=\{\frac{1}{n}: n \in \mathbb{N}\}$.  Are there any other limit points?
\end{problem}

\begin{exercise}
Provide an example of a set $S$ such that 1 is a limit point of $S$, $1\neq S$, and $S$ contains no intervals.
\end{exercise}

\begin{exercise}
Provide an example of a set $T$ with exactly two limit points.
\end{exercise}

\begin{theorem}
If $p\in\mathbb{R}$, then $p$ is a limit point of $\mathbb{Q}$.
\end{theorem}

\begin{definition}
A set is called \textbf{closed} iff it contains all of its limit points.
\end{definition}

\begin{exercise}
Provide an example of each of the following.  You do not need to prove that your answers are correct.
\begin{enumerate}
\item A closed set.
\item A set that is not closed.
\item A set that is open and closed.
\item A set that neither open or closed.
\end{enumerate}

\end{exercise}

\begin{theorem}[*]
The set $[a,b]$ is closed.
\end{theorem}

\begin{theorem}
The set $U$ is open iff $U^C$ is closed.
\end{theorem}

\begin{theorem}[*]
Every finite set is closed.
\end{theorem}

\begin{theorem}
The real numbers are both open and closed.
\end{theorem}

\begin{theorem}
The rational numbers are neither open or closed.  (You may use the fact that between any two rational numbers, there exists an irrational number.)
\end{theorem}

\begin{theorem}
The empty set is both open and closed.
\end{theorem}

\begin{theorem}[*]
Let $\{A_{\alpha}\}_{\alpha\in\Delta}$ be a collection of closed sets.  Then
\[
\bigcap_{\alpha\in \Delta} A_{\alpha}
\]
is a closed set.
\end{theorem}

\begin{exercise}
Provide an example of a collection of closed sets $\{A_{\alpha}\}_{\alpha\in\Delta}$ such that 
\[
\bigcup_{\alpha\in \Delta} A_{\alpha}
\]
is a \emph{not} closed set.
\end{exercise}

\begin{problem}
Let $A$ and $B$ be closed sets.  Determine whether $A\cup B$ is necessarily closed and prove that your answer is correct.
\end{problem}

\end{subsection}

\end{section}

\end{document}

\end{chapter}

\begin{chapter}{Induction}
%\documentclass[11pt]{article}

\usepackage{amsfonts}
\usepackage{amsmath}
\usepackage{wasysym}
\usepackage{amssymb}
\usepackage{stmaryrd}
\usepackage{amsthm}
\usepackage{fancyhdr}
\usepackage[margin=1in]{geometry}
\usepackage[hang,flushmargin,symbol*]{footmisc}
\usepackage{color}
\definecolor{darkblue}{rgb}{0, 0, .6}
\definecolor{grey}{rgb}{.7, .7, .7}
\usepackage[breaklinks]{hyperref}
\hypersetup{
	colorlinks=true,
	linkcolor=darkblue,
	anchorcolor=darkblue,
	citecolor=darkblue,
	pagecolor=darkblue,
	urlcolor=darkblue,
	pdftitle={},
	pdfauthor={}
}

\pagestyle{fancy}

\lhead{\scriptsize Course Notes for Introduction to Proof (Version 1.1)} 
\rhead{\scriptsize Instructor: \href{http://danaernst.com}{D.C. Ernst}}
\lfoot{\scriptsize This work is an adaptation of notes written by Stan Yoshinobu of Cal Poly and Matthew Jones of California State University, Dominguez Hills.}
\cfoot{}

\renewcommand{\headrulewidth}{0.4pt} 
\renewcommand{\footrulewidth}{0.4pt} 

\theoremstyle{definition}
\newtheorem{theorem}{Theorem}[section]
\newtheorem{acknowledgement}[theorem]{Acknowledgement}
\newtheorem{algorithm}[theorem]{Algorithm}
\newtheorem{axiom}[theorem]{Axiom}
\newtheorem{case}[theorem]{Case}
\newtheorem{claim}[theorem]{Claim}
\newtheorem{conclusion}[theorem]{Conclusion}
\newtheorem{condition}[theorem]{Condition}
\newtheorem{conjecture}[theorem]{Conjecture}
\newtheorem{corollary}[theorem]{Corollary}
\newtheorem{criterion}[theorem]{Criterion}
\newtheorem{definition}[theorem]{Definition}
\newtheorem{example}[theorem]{Example}
\newtheorem{exercise}[theorem]{Exercise}
\newtheorem{journal}[theorem]{Journal}
\newtheorem{lemma}[theorem]{Lemma}
\newtheorem{notation}[theorem]{Notation}
\newtheorem{problem}[theorem]{Problem}
\newtheorem{proposition}[theorem]{Proposition}
\newtheorem{remark}[theorem]{Remark}
\newtheorem{solution}[theorem]{Solution}
\newtheorem{summary}[theorem]{Summary}
\newtheorem{question}[theorem]{Question}

\begin{document}

\addtocounter{section}{2}

\begin{section}{Relations and Functions}

%\addtocounter{subsection}{3}
%\addtocounter{theorem}{58}

\begin{subsection}{Relations}

\begin{definition}
An \textbf{ordered pair} is an object of the form $(x,y)$. Two ordered pairs $(x,y)$ and $(a,b)$ are \textbf{equal} if $x=a$ and $y=b$. 
\end{definition}

\begin{definition}
An \textbf{$n$-tuple} is object of the form $(x_1, x_2,\ldots,x_n)$.  Each $x_i$ is referred to as the $i$th \textbf{component}.
\end{definition}

Note that an ordered pair is just a 2-tuple.

\begin{definition}
If $X$ and $Y$ are sets, the \textbf{Cartesian product} of $X$ and $Y$ is defined by
\[
X\times Y=\{(x,y): x\in X, y\in Y\}.
\]
That is, $X\times Y$ is the set of all ordered pairs where the first element is from $X$ and the second element is from $Y$.  The set $X\times X$ is sometimes denoted by $X^2$.  We similarly define the Cartesian product of $n$ sets, say $X_1, \ldots, X_n$, by
\[
\prod_{i=1}^{n} X_i=X_1\times \cdots \times X_n=\{(x_1,\ldots,x_n): \mbox{each } x_i\in X_i\}.
\]
\end{definition}

\begin{example}
Let $A=\{a,b,c\}$ and $B=\{\smiley,\frownie\}$.  Then 
\[
A\times B=\{(a,\smiley), (a,\frownie),(b,\smiley),(b,\frownie), (c,\smiley),(c,\frownie)\}.
\]
\end{example}

\begin{exercise}
Using the sets $A$ and $B$ from the previous example, find $B\times A$.  
\end{exercise}

\begin{exercise}
Using the set $B$ from the previous examples, find $B\times B$.  
\end{exercise}

\begin{exercise}
What general conclusion can you make about $X\times Y$ versus $Y\times X$?  When will they be equal?
\end{exercise}

\begin{exercise}
If $X$ and $Y$ are both finite sets, then how many elements will $X\times Y$ have?  Be as specific as possible.
\end{exercise}

\begin{exercise} 
Let $A=\{1, 2, 3\}$, $B=\{1,2\}$, and $C=\{1,3\}$. List the elements of the set $A \times B\times C$. 
\end{exercise}

\begin{exercise}
Let $A=\mathbb{N}$ and $B=\mathbb{R}$. Describe the elements of the set $A \times B$. 
\end{exercise}

\begin{exercise} Let $A$ be the set of all differentiable functions on the open interval $(0,1)$, and let $B$ equal the set of all derivatives of functions in $A$ evaluated at $x=\frac{1}{2}$. Describe the elements of the set $A \times B$. \end{exercise}

\begin{exercise}
Three space, $\mathbb{R}^{3}$, is a Cartesian product.  Unpack the meaning of $\mathbb{R}^{3}$ using the Cartesian product, and write the complete set notation version.
\end{exercise}

\begin{exercise}
Let $X=[0,1]$ and let $Y=\{1\}$.  Describe geometrically what $X\times Y$, $Y\times X$, $X\times X$, and $Y\times Y$ look like.
\end{exercise}

\begin{definition}
Let $X$ and $Y$ be sets. A \textbf{relation} from a set $X$ to a set $Y$ is a subset of $X \times Y$. A relation on $X$ is a subset of $X \times X$.  
\end{definition}

\begin{example}
You may not realize it, but you are familiar with many relations.  For example, on the real numbers, we have the relation $\leq$.  We could say that $(3,\pi)$ is in the relation since $3\leq \pi$.  However, $(1,-1)$ is not in the relation since $1\nleq -1$.  (Order matters!)
\end{example}

\begin{remark}
Different notations for relations are used in different contexts.  When talking about relations in the abstract, we indicate that a pair $(a,b)$ is in the relation by some notation like $a\sim b$, which is read ``$a$ is related to $b$."
\end{remark}

\begin{example}
Let $P_f$ denote the set of all people with accounts on Facebook.  Define  $F$ via $xFy$ iff $x$ is friends with $y$.  Then $F$ is a relation on $P_f$.
\end{example}

\begin{remark}
We can often represent relations using graphs or digraphs.  Given a finite set $X$ and a relation $\sim$ on $X$, a \textbf{digraph} (short for \emph{directed graph}) is a discrete graph having the members of $X$ as vertices and a directed edge from $x$ to $y$ iff $x\sim y$.
\end{remark}

\begin{example}
Let $A=\{a,b,c\}$ and define $\sim=\{(a,a),(a,b),(b,c),(c,b),(c,a)\}$.  The digraph for $\sim$ is a graph with vertices $a,b,c$ and the following arrows: $a$ to $a$, $a$ to $b$, $b$ to $c$, $c$ to $b$, $c$ to $a$.
\end{example}

\begin{exercise}
Let $A=\{1,2,3,4,5,6\}$  Define $|$ on $A$ via $x|y$ iff $x$ divides $y$.  Draw the digraph for $|$ on $A$.
\end{exercise}

When $X$ or $Y$ is infinite, it is not practical to draw a digraph.  However, you are familiar with the graphs of some relations involving infinite sets.

\begin{example}
When we write $x^2+y^2=1$, we are implicitly defining a relation.  In particular, the relation is the set of ordered pairs $(x,y)$ satisfying $x^2+y^2=1$.  In set notation:
\[
\{(x,y):x^2+y^2=1\}
\]
The graph of this relation in $\mathbb{R}^2$ is the standard unit circle.
\end{example}

\begin{exercise}
Define $\sim$ on $\mathbb{R}^2$ via $x\sim y$ iff $x\leq y$.  Draw a picture of this relation in $\mathbb{R}^2$.
\end{exercise}



\begin{definition}
Let $\sim$ be a relation on a set $A$.
\begin{enumerate}
\item $\sim$ is \textbf{reflexive} if for all $x\in A$, $x\sim x$ (every element is related to itself).
\item $\sim$ is \textbf{symmetric} if for all $x,y\in A$, if $x\sim y$, then $y\sim x$.
\item $\sim$ is \textbf{transitive} if for all $x,y,z\in A$, if $x\sim y$ and $y\sim z$, then $x\sim z$.
\end{enumerate}
\end{definition}

\begin{example}\
\begin{enumerate}
\item $\leq$ on $\mathbb{R}$ is reflexive and transitive, but not symmetric.
$<$ on $\mathbb{R}$ is transitive, but not symmetric and not reflexive.
\item If $S$ is a set, then $\subseteq$ on $\mathcal{P}(S)$ is reflexive and transitive, but not symmetric.
\item $=$ on $\mathbb{R}$ is reflexive, symmetric, and transitive.
\end{enumerate}

\end{example}

\begin{exercise}
Given a finite set $A$ and a relation $\sim$, describe what each of reflexive, symmetric, and transitive look like in terms of a digraph.
\end{exercise}

\begin{exercise}
Let $P$ be the set of people at a party and define $N$ via $(x,y)\in N$ iff $x$ knows the name of $y$.  Describe what it would mean for $N$ to be reflexive, symmetric, and transitive.
\end{exercise}

\begin{exercise}
Determine whether each of the following relations is reflexive, symmetric, or transitive.

\begin{enumerate}
\item Let $P_f$ denote the set of all people with accounts on Facebook.  Define  $F$ via $xFy$ iff $x$ is friends with $y$. 
\item Let $P$ be the set of all people and define $H$ via $xHy$ iff $x$ and $y$ have the same height.
\item Let $P$ be the set of all people and define $T$ via $xTy$ iff $x$ is taller than $y$.
\item Consider the relation ``divides" on $\mathbb{N}$.
\item Let $L$ be the set of lines and define $||$ via $l_1||l_2$ iff $l_1$ is parallel to $l_2$.
\item Let $C[0,1]$ be the set of continuous functions on $[0,1]$.  Define $f\sim g$ iff
\[
\int_0^1|f(x)|\ dx=\int_0^1|g(x)|\ dx.
\]
\item Define $\sim$ on $\mathbb{N}$ via $n\sim m$ iff $n+m$ is even.
\item Define $D$ on $\mathbb{R}$ via $(x,y)\in D$ iff $x=2y$.
\end{enumerate}
\end{exercise}

\end{subsection}

\end{section}

\end{document}
\end{chapter}

\begin{chapter}{Relations and Functions}
%\documentclass[11pt]{article}

\usepackage{amsfonts}
\usepackage{amsmath, amsthm}
\usepackage{wasysym}
\usepackage{graphicx}
\usepackage{amssymb}
\usepackage{stmaryrd}
\usepackage{amsthm}
\usepackage{fancyhdr}
\usepackage[margin=1in]{geometry}
\usepackage[hang,flushmargin,symbol*]{footmisc}
\usepackage{color}
\definecolor{darkblue}{rgb}{0, 0, .6}
\definecolor{grey}{rgb}{.7, .7, .7}
\usepackage[breaklinks]{hyperref}
\hypersetup{
	colorlinks=true,
	linkcolor=darkblue,
	anchorcolor=darkblue,
	citecolor=darkblue,
	pagecolor=darkblue,
	urlcolor=darkblue,
	pdftitle={},
	pdfauthor={}
}

\newcommand{\dom}{\operatorname{Dom}}
\newcommand{\codom}{\operatorname{Codom}}
\newcommand{\range}{\operatorname{Rng}}

\newcommand{\ndv}{\hspace{-4pt}\not|\hspace{2pt}}

\pagestyle{fancy}

\lhead{\scriptsize Notes for an Introduction to Proof Course (Version Spring 2013)}
\rhead{\scriptsize Instructor: \href{http://danaernst.com}{D.C. Ernst}}
\lfoot{\scriptsize This work is an adaptation of notes written by Stan Yoshinobu of Cal Poly and Matthew Jones of California State University, Dominguez Hills.} 
\cfoot{}
\renewcommand{\headrulewidth}{0.4pt} 
\renewcommand{\footrulewidth}{0.4pt}


\theoremstyle{definition}
\newtheorem{theorem}{Theorem}[section]
\newtheorem{acknowledgement}[theorem]{Acknowledgement}
\newtheorem{algorithm}[theorem]{Algorithm}
\newtheorem{axiom}[theorem]{Axiom}
\newtheorem{case}[theorem]{Case}
\newtheorem{claim}[theorem]{Claim}
\newtheorem{conclusion}[theorem]{Conclusion}
\newtheorem{condition}[theorem]{Condition}
\newtheorem{conjecture}[theorem]{Conjecture}
\newtheorem{corollary}[theorem]{Corollary}
\newtheorem{criterion}[theorem]{Criterion}
\newtheorem{definition}[theorem]{Definition}
\newtheorem{example}[theorem]{Example}
\newtheorem{exercise}[theorem]{Exercise}
\newtheorem{journal}[theorem]{Journal}
\newtheorem{lemma}[theorem]{Lemma}
\newtheorem{notation}[theorem]{Notation}
\newtheorem{problem}[theorem]{Problem}
\newtheorem{proposition}[theorem]{Proposition}
\newtheorem{remark}[theorem]{Remark}
\newtheorem{solution}[theorem]{Solution}
\newtheorem{summary}[theorem]{Summary}
\newtheorem{question}[theorem]{Question}
\newtheorem{skeleton}[theorem]{Skeleton Proof}

\newsavebox{\savepar}
\newenvironment{textbox}{\noindent\begin{lrbox}{\savepar}\begin{minipage}[c]{.98\textwidth}}{\end{minipage}\end{lrbox}\fcolorbox{black}{white}{\usebox{\savepar}}}

\begin{document}

\addtocounter{section}{3}

\begin{section}{Two Famous Theorems}

\begin{subsection}{The infinitude of primes}

The highlight of this section\footnote{This section is derived from work of \href{}{Dave Richeson} of Dickenson College.} is Theorem~\ref{thm:infprimes}, which states that there are infinitely primes. In case you forgot, here is the definition of a prime number.

\begin{definition}
A natural number $p$ is called \textbf{prime} iff $p$ is divisible by exactly two distinct natural numbers (namely, 1 and $p$ itself).
\end{definition}

\begin{exercise}
Is 1 a prime number?
\end{exercise}

The first known proof of Theorem~\ref{thm:infprimes} is in Eulcid's \emph{Elements} (c.\ 300 BCE). Euclid stated it as follows: 
\begin{quote}
\textbf{Proposition IX.20.} Prime numbers are more than any assigned multitude of prime numbers.
\end{quote}
There are a few interesting observations to make about Euclid's proposition and his proof. First, notice that the statement of the theorem does not contain the word ``infinity.'' The Greek's were skittish about the idea of infinity. Thus he proved that there were more primes than any given finite number. Today we'd say that they are infinite. In fact, Euclid proved that there are more than \emph{three} primes and concluded that there were more than any finite number. While we you would lose points for such a proof in this class, we can forgive Euclid for this less-than-rigorous proof;  in fact, it is easy to turn his proof into the general one that you will give below. Lastly, Euclid's proof was geometric. He was viewing his numbers as line segments with integral length. The modern concept of number was not developed yet.

Before tackling a proof of Theorem~\ref{thm:infprimes}, we need a few tools.  First, recall Problem~1.15, which asked us to prove or provide a counter example to the following statement:
\begin{quote}
Assume $a$ and $n$ are integers.  If $a$ divides $n^2$, then $a$ divides $n$.
\end{quote}
It turns out that this statement is false.  However, what if $a$ is prime?

blah blah

\begin{theorem}
\label{thm:infprimes}
There are infinitely many prime numbers.
\end{theorem}


To prove this theorem we need two lemmas. The proof of the first lemma is provided for you. 

\begin{lemma}
\label{lem:divisorsof1}
The only integers that divide $1$ are $1$ and $-1$.  
\end{lemma}
%\begin{proof}
%Let $m$ be a integer that divides $1$. Then $m$ is positive or negative.\\
%Case 1. Suppose $m>0$. We know that $m\ge 1$ because 1 is the smallest positive integer. Because $m|1$ there exists $n\in \mathbb{Z}$ such that $1=mn$. Moreover, we know that $n>0$ because $m>0$ and $mn>0$; hence $1\le n$. Multiplying both sides of $1\le n$ by $m$ gives $m\le mn=1$. Thus $1\le m\le 1$, which implies $m=1$.\\
%Case 2. Suppose $m<0$. This case is proved similarly and from it we conclude that $m=-1$.
%\end{proof}


\begin{lemma}
\label{lem:plus1}
Let $p$ be prime and $n\in \mathbb{Z}$. If $p|n$, then $p\ndv (n+1)$.
\end{lemma}

%Before we prove Theorem \ref{thm:sqrt2} we prove the following lemma.
%
%\begin{lemma}%add this to primes section! 
%\label{lem:evensquare}
%Let $n\in \mathbb{Z}$. If $n^{2}$ is even, then $n$ is even.
%\end{lemma}

\end{subsection}

\end{section}

\end{document}
%\documentclass[11pt]{article}

\usepackage{amsfonts}
\usepackage{amsmath}
\usepackage{wasysym}
\usepackage{amssymb}
\usepackage{stmaryrd}
\usepackage{amsthm}
\usepackage{fancyhdr}
\usepackage[margin=1in]{geometry}
\usepackage[hang,flushmargin,symbol*]{footmisc}
\usepackage{color}
\definecolor{darkblue}{rgb}{0, 0, .6}
\definecolor{grey}{rgb}{.7, .7, .7}
\usepackage[breaklinks]{hyperref}
\hypersetup{
	colorlinks=true,
	linkcolor=darkblue,
	anchorcolor=darkblue,
	citecolor=darkblue,
	pagecolor=darkblue,
	urlcolor=darkblue,
	pdftitle={},
	pdfauthor={}
}

\pagestyle{fancy}

\lhead{\scriptsize Course Notes for Introduction to Proof (Version 1.1)} 
\rhead{\scriptsize Instructor: \href{http://danaernst.com}{D.C. Ernst}}
\lfoot{\scriptsize This work is an adaptation of notes written by Stan Yoshinobu of Cal Poly and Matthew Jones of California State University, Dominguez Hills.}
\cfoot{}

\renewcommand{\headrulewidth}{0.4pt} 
\renewcommand{\footrulewidth}{0.4pt} 

\theoremstyle{definition}
\newtheorem{theorem}{Theorem}[section]
\newtheorem{acknowledgement}[theorem]{Acknowledgement}
\newtheorem{algorithm}[theorem]{Algorithm}
\newtheorem{axiom}[theorem]{Axiom}
\newtheorem{case}[theorem]{Case}
\newtheorem{claim}[theorem]{Claim}
\newtheorem{conclusion}[theorem]{Conclusion}
\newtheorem{condition}[theorem]{Condition}
\newtheorem{conjecture}[theorem]{Conjecture}
\newtheorem{corollary}[theorem]{Corollary}
\newtheorem{criterion}[theorem]{Criterion}
\newtheorem{definition}[theorem]{Definition}
\newtheorem{example}[theorem]{Example}
\newtheorem{exercise}[theorem]{Exercise}
\newtheorem{journal}[theorem]{Journal}
\newtheorem{lemma}[theorem]{Lemma}
\newtheorem{notation}[theorem]{Notation}
\newtheorem{problem}[theorem]{Problem}
\newtheorem{proposition}[theorem]{Proposition}
\newtheorem{remark}[theorem]{Remark}
\newtheorem{solution}[theorem]{Solution}
\newtheorem{summary}[theorem]{Summary}
\newtheorem{question}[theorem]{Question}

\begin{document}

\addtocounter{section}{3}

\begin{section}{Relations and Functions (continued)}

\addtocounter{subsection}{1}
\addtocounter{theorem}{27}

\begin{subsection}{Equivalence Relations}

\begin{remark}
So that we have them handy, let's recall the following definitions. Let $\sim$ be a relation on a set $A$.  Then
\begin{enumerate}
\item $\sim$ is \textbf{reflexive} if for all $x\in A$, $x\sim x$ (every element is related to itself).
\item $\sim$ is \textbf{symmetric} if for all $x,y\in A$, if $x\sim y$, then $y\sim x$.
\item $\sim$ is \textbf{transitive} if for all $x,y,z\in A$, if $x\sim y$ and $y\sim z$, then $x\sim z$.
\end{enumerate}
As we've seen in the previous section of notes, these conditions are mutually exclusive.  That is, a relation may have some combination of these properties, but not necessarily all of them.  However, we have a special name for when a relation does satisfy all three.
\end{remark}

\begin{definition}
Let $\sim$ be a relation on a set $A$.  Then $\sim$ is called an \textbf{equivalence relation} if $\sim$ is reflexive, symmetric, and transitive.
\end{definition}

\begin{exercise}
Given a finite set $A$ and a relation $\sim$ on $A$, describe what the corresponding digraph would have to look like in order for $\sim$ to be an equivalence relation.
\end{exercise}

\begin{exercise}\label{exer:made up}
Let $A=\{a,b,c,d,e\}$.  Make up an equivalence relation on $A$ by drawing a digraph such that $a$ is not related $b$ and $c$ is not related to $b$.
\end{exercise}

\begin{exercise}\label{exer:digraph}
Let $S=\{1,2,3,4,5,6\}$ and define
\[
\sim=\{(1,1),(1,6),(2,2),(2,3),(2,4),(3,3),(3,2),(3,4),(4,4),(4,2),(4,3),(5,5),(6,6),(6,1)\}.
\]
Justify that this is an equivalence relation.
\end{exercise}

\begin{problem}
Determine which of the following are equivalence relations.  Some of these occurred in the last section of notes and you are welcome to use your answers from those problems.

\begin{enumerate}\label{exer:lots of them}
\item\label{exer:facebook} Let $P_f$ denote the set of all people with accounts on Facebook.  Define  $F$ via $xFy$ iff $x$ is friends with $y$. 
\item Let $P$ be the set of all people and define $H$ via $xHy$ iff $x$ and $y$ have the same height.
\item Let $P$ be the set of all people and define $T$ via $xTy$ iff $x$ is taller than $y$.
\item Consider the relation ``divides" on $\mathbb{N}$.
\item Let $L$ be the set of lines and define $||$ via $l_1||l_2$ iff $l_1$ is parallel to $l_2$.
\item Let $C[0,1]$ be the set of continuous functions on $[0,1]$.  Define $f\sim g$ iff
\[
\int_0^1|f(x)|\ dx=\int_0^1|g(x)|\ dx.
\]
\item Define $\sim$ on $\mathbb{N}$ via $n\sim m$ iff $n+m$ is even.
\item Define $D$ on $\mathbb{R}$ via $(x,y)\in D$ iff $x=2y$.
\item\label{exer:mod 5} Define $\sim$ on $\mathbb{Z}$ via $a\sim b$ iff $a-b$ is a multiple of 5.
\item Define $\sim$ on $\mathbb{R}^2$ via $(x_1,y_1)\sim (x_2,y_2)$ iff $x_1^2+y_1^2=x_2^2+y_2^2$.
\item Define $\sim$ on $\mathbb{R}$ via $x\sim y$ iff $\lfloor x\rfloor =\lfloor y\rfloor$, where $\lfloor x\rfloor$ is the greatest integer less than or equal to $x$ (e.g., $\lfloor \pi\rfloor=3$, $\lfloor -1.5\rfloor=-2$, and $\lfloor 4\rfloor=4$).
\item Define $\sim$ on $\mathbb{R}$ via $x \sim y$ iff $|x-y|<1$.
\end{enumerate}
\end{problem}

\begin{definition}
Let $\sim$ be a relation on a set $A$ (not necessarily an equivalence relation) and let $x\in A$.  Then we define the \textbf{set of relatives of $x$} via
\[
R_x=\{y\in A: x\sim y\}.
\]
Also, define
\[
\Omega_{\sim}=\{R_x:x\in A\}.
\]
Notice that $\Omega_{\sim}$ is a set of sets.  In particular, an element in $\Omega_{\sim}$ is a subset of $A$.
\end{definition}

\begin{exercise}
Let $P_f$ and $F$ be as in part \ref{exer:facebook} of Exercise \ref{exer:lots of them}.  Describe $R_{\text{Bob}}$ (assume you know which Bob we're talking about).  What is $\Omega_F$?
\end{exercise}

\begin{exercise}
Using your digraph in Exercise \ref{exer:made up}, find $\Omega_{\sim}$ for all $x\in A$.  
\end{exercise}

\begin{exercise}
Consider the relation $\leq$ on $\mathbb{R}$.  If $x\in \mathbb{R}$, what is $R_x$?
\end{exercise}

\begin{exercise}
Find $R_1$ and $R_{2}$ for the relation given in part \ref{exer:mod 5} of Exercise \ref{exer:lots of them}.  How many different sets of relatives are there?  What are they?
\end{exercise}

\begin{exercise}
Find $R_x$ for all $x\in S$ for $S$ and $\sim$ from Exercise \ref{exer:digraph}.  Any observations?
\end{exercise}

\begin{theorem}[*]
Suppose $\sim$ is an equivalence relation on a set $A$ and let $a,b\in A$.  Then $R_a=R_b$ iff $a\sim b$.
\end{theorem}

\begin{theorem}[*]\label{thm:equiv yields partition}
Suppose $\sim$ is an equivalence relation on a set $A$.  Then
\begin{enumerate}
\item $\bigcup_{x\in A}R_x=A$, and
\item for all $x,y\in A$, either $R_x=R_y$ or $R_x\cap R_y=\emptyset$.
\end{enumerate}
\end{theorem}

\begin{definition}
In light of Theorem \ref{thm:equiv yields partition}, if $\sim$ is an equivalence relation on a set $A$, then we refer to each $R_x$ as the \textbf{equivalence class} of $x$.  In this case, $\Omega_{\sim}$ is the set of equivalence classes determined by $\sim$.
\end{definition}

\begin{remark}
The upshot of Theorem \ref{thm:equiv yields partition} is that given an equivalence relation, every element lives in exactly one equivalence class.  We'll see in the next section of notes that we can run this in reverse.  That is, if we separate out the elements of a set so that every element is an element of exactly one subset (like the bins of my kid's toys), then this determines an equivalence relation.  More on this later.
\end{remark}

\begin{example}
The set of relatives that you found in part \ref{exer:mod 5} of Exercise \ref{exer:lots of them} is the set of equivalence classes modulo 5.
\end{example}

\begin{exercise}
If $\sim$ is an equivalence relation on a finite set $A$, then what is the connection between the equivalence classes and the corresponding digraph?
\end{exercise}

\begin{exercise}
For each of the equivalence relations in Exercise \ref{exer:lots of them}, describe the equivalence classes as best as you can.
\end{exercise}

\end{subsection}

\end{section}

\end{document}
%\documentclass[11pt]{article}

\usepackage{amsfonts}
\usepackage{amsmath}
\usepackage{wasysym}
\usepackage{amssymb}
\usepackage{stmaryrd}
\usepackage{amsthm}
\usepackage{fancyhdr}
\usepackage[margin=1in]{geometry}
\usepackage[hang,flushmargin,symbol*]{footmisc}
\usepackage{color}
\definecolor{darkblue}{rgb}{0, 0, .6}
\definecolor{grey}{rgb}{.7, .7, .7}
\usepackage[breaklinks]{hyperref}
\hypersetup{
	colorlinks=true,
	linkcolor=darkblue,
	anchorcolor=darkblue,
	citecolor=darkblue,
	pagecolor=darkblue,
	urlcolor=darkblue,
	pdftitle={},
	pdfauthor={}
}

\pagestyle{fancy}

\lhead{\scriptsize Course Notes for Introduction to Proof (Version 1.1)} 
\rhead{\scriptsize Instructor: \href{http://danaernst.com}{D.C. Ernst}}
\lfoot{\scriptsize This work is an adaptation of notes written by Stan Yoshinobu of Cal Poly and Matthew Jones of California State University, Dominguez Hills.}
\cfoot{}

\renewcommand{\headrulewidth}{0.4pt} 
\renewcommand{\footrulewidth}{0.4pt} 

\theoremstyle{definition}
\newtheorem{theorem}{Theorem}[section]
\newtheorem{acknowledgement}[theorem]{Acknowledgement}
\newtheorem{algorithm}[theorem]{Algorithm}
\newtheorem{axiom}[theorem]{Axiom}
\newtheorem{case}[theorem]{Case}
\newtheorem{claim}[theorem]{Claim}
\newtheorem{conclusion}[theorem]{Conclusion}
\newtheorem{condition}[theorem]{Condition}
\newtheorem{conjecture}[theorem]{Conjecture}
\newtheorem{corollary}[theorem]{Corollary}
\newtheorem{criterion}[theorem]{Criterion}
\newtheorem{definition}[theorem]{Definition}
\newtheorem{example}[theorem]{Example}
\newtheorem{exercise}[theorem]{Exercise}
\newtheorem{journal}[theorem]{Journal}
\newtheorem{lemma}[theorem]{Lemma}
\newtheorem{notation}[theorem]{Notation}
\newtheorem{problem}[theorem]{Problem}
\newtheorem{proposition}[theorem]{Proposition}
\newtheorem{remark}[theorem]{Remark}
\newtheorem{solution}[theorem]{Solution}
\newtheorem{summary}[theorem]{Summary}
\newtheorem{question}[theorem]{Question}

\begin{document}

\addtocounter{section}{3}

\begin{section}{Relations and Functions (continued)}

\addtocounter{subsection}{2}
\addtocounter{theorem}{46}

\begin{subsection}{Partitions}

\begin{remark}
The upshot of Theorems 3.40 and 3.41 is that if $\sim$ is an equivalence relation on a set $A$, then $\sim$ breaks $A$ up into pairwise disjoint chunks, where each chunk is some $R_{a}$ for $a\in A$. Furthermore, each  pair of elements in the same set of relatives are related via $\sim$.
\end{remark}

As you've probably already noticed, equivalence relations are intimately related to the following concept.

\begin{definition}
A collection $\Omega$ of nonempty subsets of a set $A$ is said to be a \textbf{partition} of $A$ if the elements of $\Omega$ satisfy:
\begin{enumerate}
\item Given $X,Y\in\Omega$, either $X=Y$ or $X\cap Y=\emptyset$ (We can't have both at the same time. Do you see why?), and
\item $\displaystyle \bigcup_{X\in\Omega}X=A$.
\end{enumerate}
That is, the elements of $\Omega$ are pairwise disjoint and their union is all of $A$.
\end{definition}

\begin{example}
The following are all examples of partitions of the given set.  Perhaps you can find exceptions in these examples, but please take them at face value.
\begin{enumerate}
\item men, women (set of people)
\item Democrat, Republican, Independent, Green Party, Libertarian, etc. (set of registered voters)
\item freshman, sophomore, junior, senior (set of high school students)
\item evens, odds (set of integers)
\item rationals, irrationals (set of real numbers)
\end{enumerate}
\end{example}

\begin{example}\label{ex:a partition}
Let $A=\{a,b,c,d,e,f\}$ and $\Omega=\{X_{1},X_{2},X_{3}\}$, where $X_{1}=\{a\}$, $X_{2}=\{b,c,d\}$, and $X_{3}=\{e,f\}$.  Then $\Omega$ is a partition of $A$ since the elements of $\Omega$ are pairwise disjoint and their union is all of $A$.
\end{example}

\begin{exercise}
Consider the set $A$ from Example \ref{ex:a partition}.
\begin{enumerate}
\item Find a partition of $A$ that has 4 subsets in the partition.
\item Find a collection of subsets of $A$ that does \emph{not} form a partition.
\end{enumerate}
\end{exercise}

\begin{exercise}
Find a partition of $\mathbb{N}$ that consists of 3 subsets, where one of the sets is finite and the remaining two sets are infinite.
\end{exercise}

\begin{exercise}
Let $P$ be the set of prime numbers, $N$ be the set of odd natural numbers that are not prime, and $E$ be the set of even natural numbers.  Explain why this is not a partition of $\mathbb{N}$.
\end{exercise}

The next theorem spells out half of the close connection between partitions and equivalence relations.  Hopefully you were anticipating this.

\begin{theorem}[*]
Let $\sim$ be an equivalence relation on a set $A$.  Then $\Omega_{\sim}$ forms a partition of $A$.
\end{theorem}

\begin{exercise}
Consider the equivalence relation
\[
\sim=\{(1,1),(1,2),(2,1), (2,2),(3,3),(4,4),(4,5),(5,4),(6,6),(5,6),(6,5),(4,6),(6,4)\}
\]
on the set $A=\{1,2,3,4,5,6\}$.  Find the partition determined by $\Omega_{\sim}$.
\end{exercise}

It turns out that we can reverse the situation, as well.  That is, given a partition, we can form an equivalence relation.  Before proving this, we need a definition.

\begin{definition}
Let $A$ be a set and $\Omega$ any collection of subsets from $\mathcal{P}(A)$ (not necessarily a partition).  If $a,b\in A$, we will define $a$ to be $\Omega$-related to $b$ if there exists an $R\in \Omega$ that contains both $a$ and $b$.  This relation is denoted by $\sim_{\Omega}$ and is called the \textbf{relation on $A$ associated to $\Omega$}.
\end{definition}

\begin{remark}
This definition may look more awkward than the actual underlying concept.  The idea is that if two elements are in the same subset, then they are related.  For example, when my kids pick up all their toys and put them in the appropriate toy bins, we say that two toys are related if they are in the same bin.
\end{remark}

\begin{remark}
Notice that we have two notations that looks similar: $\Omega_{\sim}$ and $\sim_{\Omega}$.  
\begin{enumerate}
\item $\Omega_{\sim}$ is the collection of subsets of $A$ determined by the relation $\sim$.
\item $\sim_{\Omega}$ is the relation determined by the collection of subsets $\Omega$.
\end{enumerate}
\end{remark}

\begin{exercise}
Let $A=\{a,b,c,d,e,f\}$ and let $\Omega=\{X_{1},X_{2},X_{3}\}$, where $X_{1}=\{a,c\}$, $X_{2}=\{b,c\}$, and $X_{3}=\{d,f\}$.  List the elements of $\sim_{\Omega}$ by listing ordered pairs or drawing a digraph.
\end{exercise}

\begin{exercise}
Let $A$ and $\Omega$ be as in Example \ref{ex:a partition}.  List the elements of $\sim_{\Omega}$ by listing ordered pairs or drawing a digraph.
\end{exercise}

\begin{theorem}[*]
Let $A$ be a set and let $\Omega$ be a collection of subsets from $\mathcal{P}(A)$ (not necessarily a partition).  Then $\sim_{\Omega}$ is symmetric.
\end{theorem}

\begin{exercise}
Give an example of a set $A$ and a collection $\Omega$ from $\mathcal{P}(A)$ such that the relation $\sim_{\Omega}$ is not reflexive.
\end{exercise}

\begin{theorem}[*]
Let $A$ be a set and let $\Omega$ be a collection of subsets from $\mathcal{P}(A)$.  If $\displaystyle \bigcup_{R\in\Omega}R=A$, then $\sim_{\Omega}$ is reflexive.
\end{theorem}

\begin{theorem}[*]
Let $A$ be a set and let $\Omega$ be a collection of subsets from $\mathcal{P}(A)$.  If the elements of $\Omega$ are pairwise disjoint, then $\sim_{\Omega}$ is transitive.
\end{theorem}

\begin{corollary}[*]
Let $A$ be a set and let $\Omega$ be a partition of $A$.  Then $\sim_{\Omega}$ is an equivalence relation.
\end{corollary}

\begin{remark}
The previous corollary says that every partition determines a natural equivalence relation.  Namely, two elements are related if they are in the same equivalence class.
\end{remark}

\begin{exercise}
Let $A=\{\circ, \triangle, \blacktriangle, \square, \blacksquare, \bigstar, \smiley, \frownie\}$.  Make up a partition $\Omega$ on $A$ and then draw the digraph corresponding to $\sim_{\Omega}$.
\end{exercise}

\end{subsection}

\end{section}

\end{document}
%\documentclass[11pt]{article}

\usepackage{amsfonts}
\usepackage{amsmath}
\usepackage{wasysym}
\usepackage{amssymb}
\usepackage{stmaryrd}
\usepackage{amsthm}
\usepackage{fancyhdr}
\usepackage[margin=1in]{geometry}
\usepackage[hang,flushmargin,symbol*]{footmisc}
\usepackage{color}
\definecolor{darkblue}{rgb}{0, 0, .6}
\definecolor{grey}{rgb}{.7, .7, .7}
\usepackage[breaklinks]{hyperref}
\hypersetup{
	colorlinks=true,
	linkcolor=darkblue,
	anchorcolor=darkblue,
	citecolor=darkblue,
	pagecolor=darkblue,
	urlcolor=darkblue,
	pdftitle={},
	pdfauthor={}
}

\newcommand{\dom}{\operatorname{Dom}}
\newcommand{\codom}{\operatorname{Codom}}
\newcommand{\range}{\operatorname{Rng}}

\pagestyle{fancy}

\lhead{\scriptsize Course Notes for Introduction to Proof (Version 1.1)} 
\rhead{\scriptsize Instructor: \href{http://danaernst.com}{D.C. Ernst}}
\lfoot{\scriptsize This work is an adaptation of notes written by Stan Yoshinobu of Cal Poly and Matthew Jones of California State University, Dominguez Hills.}
\cfoot{}

\renewcommand{\headrulewidth}{0.4pt} 
\renewcommand{\footrulewidth}{0.4pt} 

\theoremstyle{definition}
\newtheorem{theorem}{Theorem}[section]
\newtheorem{acknowledgement}[theorem]{Acknowledgement}
\newtheorem{algorithm}[theorem]{Algorithm}
\newtheorem{axiom}[theorem]{Axiom}
\newtheorem{case}[theorem]{Case}
\newtheorem{claim}[theorem]{Claim}
\newtheorem{conclusion}[theorem]{Conclusion}
\newtheorem{condition}[theorem]{Condition}
\newtheorem{conjecture}[theorem]{Conjecture}
\newtheorem{corollary}[theorem]{Corollary}
\newtheorem{criterion}[theorem]{Criterion}
\newtheorem{definition}[theorem]{Definition}
\newtheorem{example}[theorem]{Example}
\newtheorem{exercise}[theorem]{Exercise}
\newtheorem{journal}[theorem]{Journal}
\newtheorem{lemma}[theorem]{Lemma}
\newtheorem{notation}[theorem]{Notation}
\newtheorem{problem}[theorem]{Problem}
\newtheorem{proposition}[theorem]{Proposition}
\newtheorem{remark}[theorem]{Remark}
\newtheorem{solution}[theorem]{Solution}
\newtheorem{summary}[theorem]{Summary}
\newtheorem{question}[theorem]{Question}

\begin{document}

\addtocounter{section}{2}

\begin{section}{Relations and Functions}

\addtocounter{subsection}{3}
\addtocounter{theorem}{67}

\begin{subsection}{Introduction to Functions}

The concept of function is one of the most important and fundamental ones in the field of mathematics.  Functions are used in all branches of mathematics to model diverse situations and pull together ideas that at first seem unrelated.  Functions are as vital as number.

Undoubtably, you have encountered the concept of function in your prior mathematical experience.  In this section, we will introduce the concept of function as a special type of relation.  As you shall see, this agrees with any previous definition of function that you may have learned.  

Up until this point, you've probably only encountered functions as an algebraic rule, e.g. $f(x)=x^{2}-1$, for transforming one real number into another.  However, we can study functions in a much broader context.  Loosely speaking, the basic building blocks of a function are a first set and a second sets, say $X$ and $Y$, respectively, and a ``correspondence'' that assigns to each element of $X$ to exactly one element of $Y$.  Let's take a look at the actual definition.

\begin{definition}
Let $X$ and $Y$ be two nonempty sets.  A \textbf{function} from set $X$ to set $Y$, denoted $f:X\to Y$, is a relation (i.e., subset of $X\times Y$) such that:

\begin{enumerate}\label{def:function}
\item For each $x\in X$, there exists $y\in Y$ such that $(x,y)\in f$, and
\item If $(x,y_{1}), (x,y_{2}) \in f$, then $y_{1}=y_{2}$.
\end{enumerate}
Note that if $(x,y)\in f$, we usually write $y=f(x)$ and say that ``$f$ maps $x$ to $y$.''
\end{definition}

\begin{remark}
Item 1 of Definition \ref{def:function} says that every element of $X$ appears in the first coordinate of an ordered pair in the relation.  Item 2 says that each element of $X$ only appears once in the first coordinate of an ordered pair in the relation.  It is important to note that there are no restrictions on whether an element of $Y$ ever appears in the second coordinate.  Furthermore, if an element of $B$ appears in the second coordinate, it may appear again in a different ordered pair.
\end{remark}

\begin{definition}
The set $X$ from Definition \ref{def:function} is called the \textbf{domain} of $f$ and is denoted by $\dom(f)$.  The set $Y$ is called the \textbf{codomain} of $f$ and is denoted by $\codom(f)$.  The set
\[
\range(f)=\{y\in Y: \mbox{there exists }x\mbox{ such that } y=f(x)\}
\]
is called the \textbf{range} of $f$ or the \textbf{image of $X$} under $f$.
\end{definition}

\begin{remark}
It follows immediately from the definition that $\range(f)\subseteq \codom(f)$.  However, it is possible that the range of $f$ is strictly smaller.
\end{remark}

\begin{remark}
If $f$ is a function and $(x,y)\in f$, then we may refer to $x$ as the \textbf{input} of $f$ and $y$ as the \textbf{output} of $f$.
\end{remark}

\begin{exercise}\label{exer:lots}
Let $X=\{\circ, \square,\triangle,\smiley,\}$ and $Y=\{a,b,c,d,e\}$.  Determine whether each of the following represent functions.  Explain.  If the relation is a function, determine the domain, codomain, and range.

\begin{enumerate}
\item $f:X\to Y$ defined via $f=\{(\circ, a),(\square,b),(\triangle,c),(\smiley,d)\}$.
\item $g:X\to Y$ defined via $g=\{(\circ, a),(\square,b),(\triangle,c),(\smiley,c)\}$.
\item $h:X\to Y$ defined via $h=\{(\circ, a),(\square,b),(\triangle,c),(\circ,d)\}$.
\item $k:X\to Y$ defined via $k=\{(\circ, a),(\square,b),(\triangle,c),(\smiley,d),(\square,e)\}$.
\item $l:X\to Y$ defined via $l=\{(\circ, e),(\square,e),(\triangle,e),(\smiley,e)\}$.
\item $m:X\to Y$ defined via $m=\{(\circ, a),(\triangle,b),(\smiley,c)\}$.
\item $\operatorname{happy}:Y\to X$ defined via $\operatorname{happy}(y)=\smiley$ for all $y\in Y$.
\item $\operatorname{id}:X\to X$ defined via $\operatorname{id}(x)=x$ for all $x\in X$.
\item $\operatorname{nugget}:X\to X$ defined via 
\[
\operatorname{nugget}(x)=\begin{cases}
x, & \mbox{if } x\mbox{ is a geometric shape},\\
\square, & \mbox{otherwise}.
\end{cases}
\]
\end{enumerate}
\end{exercise}

\begin{definition}
One useful representation of functions on finite sets is via \textbf{bubble diagrams}.  To draw a bubble diagram for a function $f:X\to Y$, draw one circle (i.e, a ``bubble'') for each of $X$ and $Y$ and for each element of each set, put a dot in the corresponding set.  Typically, we draw $X$ on the left and $Y$ on the right.  Now, draw an arrow from $x\in X$ to $y\in Y$ if $f(x)=y$ (i.e., $(x,y)\in f$).  In fact, we can draw bubble diagrams even if $f$ isn't a function.
\end{definition}

\begin{exercise}
For each of the relations in Exercise \ref{exer:lots} draw the corresponding bubble diagram.
\end{exercise}

\begin{problem}
What properties does a bubble diagram have to have in order to represent a function?
\end{problem}

\begin{exercise}
Provide an example of each of the following.  You may draw a bubble diagram, write down a list of ordered pairs, or a write a formula (as long as the domain and codomain are clear).
\begin{enumerate}
\item A function $f$ from a set with 4 elements to a set with 3 elements such that $\range(f)=\codom(f)$.
\item A function $g$ from a set with 4 elements to a set with 3 elements such that $\range(g)$ is strictly smaller than $\codom(g)$.
\end{enumerate}
\end{exercise}

\begin{problem}
Let $f:X\to Y$ be a function and suppose that $X$ and $Y$ have $n$ and $m$ elements in them, respectively.  Also, suppose that $n<m$.  Is it possible for $\range(f)=\codom(f)$?  Explain.
\end{problem}

\begin{problem}
In high school I am sure that you were told that a graph represents a function if it passes the \textbf{vertical line test}.  Using our terminology of ordered pairs, explain why this works.
\end{problem}

\begin{definition}
Two functions are equal if they have the same domain, same codomain, and the same set of ordered pairs in the relation.
\end{definition}

\begin{remark}
If two functions are defined by the same algebraic formula, but have different domains, then they are \emph{not} equal.  For example, the function $f:\mathbb{R}\to \mathbb{R}$ defined via $f(x)=x^{2}$ is not equal to the function $g:\mathbb{N}\to\mathbb{N}$ defined via $g(x)=x^{2}$.
\end{remark}

\begin{theorem}
If $f:X\to Y$ and $g:X\to Y$ are functions, then $f=g$ iff $f(x)=g(x)$ for all $x\in X$.
\end{theorem}

\begin{definition}
Let $f:X\to Y$ be a function.
\begin{enumerate}
\item The function $f$ is said to be \textbf{one-to-one} (or \textbf{injective}) if for all $y\in \range(f)$, there is a unique $x\in X$ such that $y=f(x)$.
\item The function $f$ is said to be \textbf{onto} (or \textbf{surjective}) if for all $y\in Y$, there exists $x\in X$ such that $y=f(x)$.
\item If $f$ is both one-to-one and onto, we say that $f$ is a \textbf{one-to-one correspondence} (or a \textbf{bijection}).
\end{enumerate}
\end{definition}

\begin{exercise}
Provide an example of each of the following.  You may draw a bubble diagram, write down a list of ordered pairs, or write a formula (as long as the domain and codomain are clear).  Assume that $X$ and $Y$ are finite sets.
\begin{enumerate}
\item A function $f:X\to Y$ that is one-to-one but not onto.
\item A function $f:X\to Y$ that is onto but not one-to-one.
\item A function $f:X\to Y$ that is both one-to-one and onto.
\item A function $f:X\to Y$ that is neither one-to-one nor onto.
\end{enumerate}

\end{exercise}

\begin{problem}
Perhaps you've heard of the \textbf{horizontal line test} (i.e., every horizontal line hits the graph of $f:\mathbb{R}\to\mathbb{R}$ at most once).  What is the horizontal line test testing for?  Explain.
\end{problem}

\begin{exercise}
Provide an example of each of the following.  You may either draw a graph or write down a formula.  Make sure you have the correct domain.
\begin{enumerate}
\item A function $f:\mathbb{R}\to \mathbb{R}$ that is one-to-one but not onto.
\item A function $f:\mathbb{R}\to \mathbb{R}$ that is onto but not one-to-one.
\item A function $f:\mathbb{R}\to \mathbb{R}$ that is both one-to-one and onto.
\item A function $f:\mathbb{R}\to \mathbb{R}$ that is neither one-to-one nor onto.
\end{enumerate}

\end{exercise}

\begin{theorem}[*]
Let $f:X\to Y$ be a function.  Then $f$ is one-to-one iff for all $x_{1}, x_{2}\in X$, if $f(x_{1})=f(x_{2})$, then $x_{1}=x_{2}$.
\end{theorem}

\begin{remark}
The previous theorem gives a technique for proving that a given function is one-to-one.  Start by assuming that $f(x_{1})=f(x_{2})$ and then work to show that $x_{1}=x_{2}$.
\end{remark}

\begin{remark}
To show that a given function is onto, you should start with an arbitrary $y\in \range(f)$ and then work to show that there exists $x\in X$ such that $y=f(x)$.
\end{remark}

\begin{exercise}
Determine which of the following functions are one-to-one, onto, both, or neither.  In each case, you should provide proofs and counterexamples as appropriate.

\begin{enumerate}
\item $f:\mathbb{R}\to \mathbb{R}$ defined via $f(x)=x^{2}$
\item $g:\mathbb{R}\to [0,\infty)$ defined via $g(x)=x^{2}$
\item $h:\mathbb{R}\to \mathbb{R}$ defined via $g(x)=x^{3}$
\item $k:\mathbb{R}\to \mathbb{R}$ defined via $k(x)=x^{3}-x$
\item $l: \mathbb{R}\times \mathbb{R}\to \mathbb{R}$ defined via $f(x_{1},x_{2})=x_{1}^{2}+x_{2}^{2}$
\item $N:\mathbb{N}\to \mathbb{N}\times \mathbb{N}$ defined via $N(n)=(n,n)$
\end{enumerate}
\end{exercise}

\begin{exercise}
Let $A$ and $B$ be sets and let $S\subseteq A\times B$.  Define $\pi_{1}:S\to A$ and $\pi_{2}:S\to B$ via $\pi_{1}(a,b)=a$ and $\pi_{2}(a,b)=b$.  We call $\pi_{1}$ (respectively, $\pi_{2}$) the \textbf{projections} of $S$ onto $A$ (respectively $B$).
\begin{enumerate}
\item Provide examples to show that $\pi_{1}$ does not need to be one-to-one or onto.
\item Suppose that $S$ is a function (recall that a function is a set of ordered pairs, so this makes sense).  Is $\pi_{1}$ one-to-one? Is $\pi_{1}$ onto?  How about $\pi_{2}$?
\end{enumerate} 
\end{exercise}


\end{subsection}

\end{section}

\end{document}
%\documentclass[11pt]{article}

\usepackage{amsfonts}
\usepackage{amsmath}
\usepackage{wasysym}
\usepackage{amssymb}
\usepackage{stmaryrd}
\usepackage{amsthm}
\usepackage{simpsons}
\usepackage{fancyhdr}
\usepackage[margin=1in]{geometry}
\usepackage[hang,flushmargin,symbol*]{footmisc}
\usepackage{color}
\definecolor{darkblue}{rgb}{0, 0, .6}
\definecolor{grey}{rgb}{.7, .7, .7}
\usepackage[breaklinks]{hyperref}
\hypersetup{
	colorlinks=true,
	linkcolor=darkblue,
	anchorcolor=darkblue,
	citecolor=darkblue,
	pagecolor=darkblue,
	urlcolor=darkblue,
	pdftitle={},
	pdfauthor={}
}

\newcommand{\dom}{\operatorname{Dom}}
\newcommand{\codom}{\operatorname{Codom}}
\newcommand{\range}{\operatorname{Rng}}

\pagestyle{fancy}

\lhead{\scriptsize Course Notes for Logic, Proof, \& Axiomatic Systems (Spring 2011)} 
\rhead{\scriptsize Instructor: \href{http://oz.plymouth.edu/~dcernst}{D.C. Ernst}} 
\lfoot{\scriptsize This work is an adaptation of notes written by Stan Yoshinobu of Cal Poly and Matthew Jones of California State University, Dominguez Hills.} 
\cfoot{} 
\renewcommand{\headrulewidth}{0.4pt} 
\renewcommand{\footrulewidth}{0.4pt} 

\theoremstyle{definition}
\newtheorem{theorem}{Theorem}[section]
\newtheorem{acknowledgement}[theorem]{Acknowledgement}
\newtheorem{algorithm}[theorem]{Algorithm}
\newtheorem{axiom}[theorem]{Axiom}
\newtheorem{case}[theorem]{Case}
\newtheorem{claim}[theorem]{Claim}
\newtheorem{conclusion}[theorem]{Conclusion}
\newtheorem{condition}[theorem]{Condition}
\newtheorem{conjecture}[theorem]{Conjecture}
\newtheorem{corollary}[theorem]{Corollary}
\newtheorem{criterion}[theorem]{Criterion}
\newtheorem{definition}[theorem]{Definition}
\newtheorem{example}[theorem]{Example}
\newtheorem{exercise}[theorem]{Exercise}
\newtheorem{journal}[theorem]{Journal}
\newtheorem{lemma}[theorem]{Lemma}
\newtheorem{notation}[theorem]{Notation}
\newtheorem{problem}[theorem]{Problem}
\newtheorem{proposition}[theorem]{Proposition}
\newtheorem{remark}[theorem]{Remark}
\newtheorem{solution}[theorem]{Solution}
\newtheorem{summary}[theorem]{Summary}
\newtheorem{question}[theorem]{Question}

\begin{document}

\addtocounter{section}{2}

\begin{section}{Relations and Functions}

\addtocounter{subsection}{4}
\addtocounter{theorem}{90}

\begin{subsection}{Compositions and Inverses}

\begin{definition}
If $f:X\to Y$ and $g:Y\to Z$ are functions, then a new function $g\circ f:X\to Z$ can be defined by $(g\circ f)(x)=g(f(x))$ for all $x\in\dom(f)$.
\end{definition}

\begin{remark}
It is important to notice that the function on the right is the one that ``goes first.''
\end{remark}

\begin{exercise}
In each case, give examples of finite sets $X$, $Y$, and $Z$, and functions $f:X\to Y$ and $g:Y\to Z$ that satisfy the given conditions.  Drawing bubble diagrams is sufficient.
\begin{enumerate}
\item $f$ is onto, but $g\circ f$ is not onto.
\item $g$ is onto, but $g\circ f$ is not onto.
\item $f$ is one-to-one, but $g\circ f$ is not one-to-one.
\item $g$ is one-to-one, but $g\circ f$ is not.
\end{enumerate}
\end{exercise}

\begin{theorem}[*]
If $f:X\to Y$ and $g:Y\to Z$ are both functions that are onto, then $g\circ f$ is also onto.
\end{theorem}

\begin{theorem}[*]
If $f:X\to Y$ and $g:Y\to Z$ are both functions that are one-to-one, then $g\circ f$ is also one-to-one.
\end{theorem}

\begin{corollary}
If $f:X\to Y$ and $g:Y\to Z$ are both one-to-one correspondences, then $g\circ f$ is also a one-to-one correspondence.
\end{corollary}

\begin{problem}
Assume that $f:X\to Y$ and $g:Y\to Z$ are both functions.  For each of the following statements, if the statement is true, then prove it.  If the statement is false, provide a counterexample.
\begin{enumerate}
\item If $g\circ f$ is one-to-one, then $f$ is one-to-one.
\item If $g\circ f$ is one-to-one, then $g$ is one-to-one.
\item If $g\circ f$ is onto, then $f$ is onto.
\item If $g\circ f$ is onto, then $g$ is onto.
\end{enumerate}
\end{problem}

\begin{definition}
Let $f:X\to Y$ be a function.  The relation $f^{-1}$, called \textbf{$f$ inverse}, is defined via
\[
f^{-1}=\{(f(x),x):x\in X\}.
\]
\end{definition}

\begin{remark}
Notice that we called $f^{-1}$ a relation and not a function.  In some circumstances $f^{-1}$ will be a function and sometimes it won't be.
\end{remark}

\begin{exercise}
Provide an example of a function $f:X\to Y$ such that $f^{-1}$ is \emph{not} a function.  A bubble diagram is sufficient.
\end{exercise}

\begin{exercise}
Provide an example of a function $f:X\to Y$ such that $f^{-1}$ is a function. A bubble diagram is sufficient.
\end{exercise}

\begin{theorem}[*]
Let $f:X\to Y$ be a function.  Then $f^{-1}$ is a function iff $f$ is \underline{\ \ \ \ \ \ \ \ \ \ \ \ \ \ \ \ \ \ \ \ \ \ \ \ }.
\end{theorem}

\begin{theorem}[*]\label{thm:comp of inverses}
Let $f:X\to Y$ be a function and suppose that $f^{-1}$ is a function.  Then
\begin{enumerate}
\item $(f\circ f^{-1})(x)=x$ for all $x\in Y$, and
\item $(f^{-1}\circ f)(x)=x$ for all $x\in X$.
\end{enumerate}
(You only need to prove one of these statements; the other is similar.)
\end{theorem}

\begin{theorem}[*]\label{thm:unique inverse}
Let $f:X\to Y$ and $g:Y\to X$ be functions such that $f$ is a one-to-one correspondence.  If $(f\circ g)(x)=x$ for all $x\in Y$ and $(g\circ f)(x)=x$ for all $x\in X$, then $g=f^{-1}$.
\end{theorem}

\begin{remark}
The upshot of the previous two theorems is that if $f^{-1}$ is a function, then it is the only one satisfying the two-sided ``undoing'' property exhibited in Theorem \ref{thm:comp of inverses}.
\end{remark}

The next theorem can be considered to be a converse of Theorem \ref{thm:unique inverse}.

\begin{theorem}[*]
Let $f:X\to Y$ and $g:Y\to X$ be functions satisfying $(f\circ g)(x)=x$ for all $x\in Y$ and $(g\circ f)(x)=x$ for all $x\in X$.  Then $f$ is a one-to-one correspondence.
\end{theorem}

\begin{theorem}[*]
Let $f:X\to Y$ and $g:Y\to Z$ be functions.  If $f$ and $g$ are both one-to-one correspondences, then $(g\circ f)^{-1}=f^{-1}\circ g^{-1}$.
\end{theorem}

\end{subsection}

\end{section}

\end{document}
\end{chapter}
\end{comment}

\end{document}

%%% Local Variables:
%%% mode: latex
%%% TeX-master: t
%%% End:
